\documentclass[11pt,a4paper]{article}

\usepackage[utf8]{inputenc}
\usepackage[T1]{fontenc}
\usepackage{amsmath,amssymb,amsthm}
\usepackage{mathtools}
\usepackage{hyperref}
\usepackage[margin=1in]{geometry}
\usepackage{enumitem}
\usepackage{microtype}
\usepackage{array}
\usepackage{booktabs}
\usepackage{url}
\usepackage{etoolbox}
\usepackage{longtable}
\usepackage{needspace}
\usepackage[most]{tcolorbox}

% Align with suite §2.8 overfull policy (see NEMS suite_preamble).
\setlength{\hfuzz}{3pt}

\newenvironment{keybox}[1]{%
  \Needspace*{8\baselineskip}%
  \begin{tcolorbox}[
    enhanced,
    breakable=false,
    colback=blue!5,
    colframe=blue!45!black,
    coltitle=black,
    colbacktitle=blue!14,
    boxrule=0.85pt,
    arc=2.5pt,
    left=9pt, right=9pt, top=6pt, bottom=6pt,
    titlerule=0.55pt,
    fonttitle=\bfseries,
    title={#1},
    before upper=\raggedright,
    before skip=0.9\baselineskip,
    after skip=0.9\baselineskip
  ]%
}{%
  \end{tcolorbox}%
}

% Theorem environments
\theoremstyle{plain}
\newtheorem{theorem}{Theorem}[section]
\newtheorem{proposition}[theorem]{Proposition}
\newtheorem{lemma}[theorem]{Lemma}
\newtheorem{corollary}[theorem]{Corollary}

\theoremstyle{definition}
\newtheorem{definition}[theorem]{Definition}
\newtheorem{example}[theorem]{Example}

\theoremstyle{remark}
\newtheorem{remark}[theorem]{Remark}
\newtheorem{notation}[theorem]{Notation}

% Allow overfull in remark environments and longtable (Lean names with underscores don't break)
\AtBeginEnvironment{remark}{\small\sloppy}
\AtBeginEnvironment{longtable}{\sloppy}

% Custom commands
\newcommand{\Bare}{\mathsf{Bare}}
\newcommand{\Real}{\mathsf{Realized}}
\newcommand{\compare}{\pi}
\newcommand{\Fib}{\mathrm{Fib}}
\newcommand{\Ker}{\mathcal{K}}
\newcommand{\KerAt}{\mathcal{K}}
\newcommand{\BD}{\mathrm{BD}}
\newcommand{\RCS}{\mathsf{RCS}}

\title{\textbf{The Geometry of What Maps Forget}\\[6pt]
\large Residual kernels, observable duality, and exact resolution complexity}

\author{Nova Spivack\\[0.3em]
{\small Independent Researcher}}
\date{March 2026}

\begin{document}

\maketitle

\noindent\emph{This paper develops the universal residual geometry that appears once comparison ceases to be injective.}
\medskip

\noindent\textbf{Keywords:} residual kernel, information loss, fiber automorphism,
observable duality, resolution complexity, graph coloring, Lean~4 formalization.

\noindent\textbf{MSC 2020:} 03B70, 03F40, 18A05, 05C15.

\noindent\textbf{Code availability:} Lean~4 formalization in
\texttt{reflexive-architecture-lean} (main residual/universal development) and,
for the supplementary bridge discussed in \S\ref{sec:representation-as-structured-forgetting},
\texttt{reflexive-closure-lean}. Build (from the corresponding repository): \texttt{lake build}.
The lemmas this paper cites as machine-checked are \texttt{sorry}-free in those sources.
Appendix~\ref{app:index} matches named paper results to Lean declaration names and file paths;
Remark~\ref{rem:structured-forgetting-lean} lists the adjoining bridge modules.

\begin{abstract}
\raggedright
We prove that whenever a comparison map fails to be injective, the resulting residual is not vague leftover structure but a precise geometric object: a kernel decomposed fiberwise, with exact resolution cost, monotone refinement behavior, and functorial invariants. We then show that this residual geometry is universal: every comparison map carries it, and adequacy determines when it is nontrivial.

For an arbitrary map $\compare : \Real \to \Bare$, five conditions---non-exhaustion,
non-injectivity, existence of a non-trivial $\compare$-preserving
endomorphism, witness diversity, and nonemptiness of the residual
kernel---are proved equivalent. Either all hold simultaneously, or none
do (all-or-nothing dichotomy).

We identify the \emph{residual kernel} $\Ker(\compare)$ as the canonical
object of information loss: the set of distinct domain pairs sharing a
common image. We prove its complete geometry: it decomposes
fiberwise into complete indistinguishability cliques; bare-determined
observables form a Boolean algebra dual to kernel separation; and
refinement acts as graph coloring of fiber cliques, with exact per-fiber
complexity (resolving a fiber of size~$n$ requires exactly~$n$ colors).
The comparison map is the universal coarse observable; the structure is
functorial; and the entire theory holds unconditionally for every function
between any two types.

Because the theory is stated for an arbitrary function between arbitrary
types, it applies not only to certification but to any map that compresses
structure: group homomorphisms, topological quotients, type-theoretic
projections, measurements, compilers, and hash functions all inherit the
full residual geometry. The theory is a universal geometry of information
loss under any map.

All results are formalized in Lean~4~\cite{Lean4} with Mathlib~\cite{Mathlib}, with zero
\texttt{sorry} in proof terms on the formalization path summarized in Appendix~\ref{app:index}. The theory is discharged against the
Infinity Compression flagship and admits instances from algebra, topology,
and computability.

Section~\ref{sec:representation-as-structured-forgetting} then develops an interpretive synthesis---representation as selective forgetting, the ascent from map to function to ontology, and reflexive non-exhaustion---explicitly layered on the formal spine (Remark~\ref{rem:representation-interpretive}).

\textbf{Trust boundary:} Universal claims are theorems in Lean for arbitrary types and comparison maps in the \texttt{ReflexiveArchitecture.Universal} hierarchy (Appendix~\ref{app:index}); adequacy hypotheses are explicit where enrichment is forced. \emph{Reflexive-closure-lean} supplies a small additional bridge (map forgetting, representation/remainder interfaces, and class-level invariants) named in Remark~\ref{rem:structured-forgetting-lean}; the physical reading of conservation as ``law'' remains interpretive (\S\ref{sec:representation-conservation}). Philosophical claims elsewhere in \S\ref{sec:representation-as-structured-forgetting} are interpretive unless noted as restating the formal fiber geometry or those abstract interfaces.
\end{abstract}

\section{Premises, disagreement coordinates, and trust boundary}\label{sec:premises-geometry-maps-forget}
\begin{keybox}{Where disagreement can (and cannot) live}
\textbf{Formal spine.}
The five-way equivalence and residual-kernel geometry are proved in the cited \texttt{reflexive-architecture-lean} modules; injective countercases are part of the stated dichotomy.

\medskip
\textbf{Premise coordinates (for external readers).}
{\raggedright
\begin{itemize}[leftmargin=2em]
\item \textbf{Certification rhetoric vs.\ minimal hypotheses.} Examples motivate bare/enriched language; theorems assume only a map \(\compare : \Real \to \Bare\) unless a section invokes adequacy structure.
  \emph{Logical effect:} Rejecting certification metaphysics does not touch the bare type-theoretic statements.
\item \textbf{Summit route vs.\ partial formalizations.} The Lean scope for this paper is the residual/\texttt{Universal} route mirrored in Appendix~\ref{app:index}; partial homotopy-layer gaps in other branches do not retroactively weaken that route if not claimed here.
  \emph{Logical effect:} Audit claims against Appendix~\ref{app:index} and the cited \texttt{.lean} paths.
\end{itemize}
\par}
\end{keybox}

\tableofcontents

%% ========================================================================
\section{Introduction}\label{sec:intro}
%% ========================================================================

Fix a comparison map $\compare : \Real \to \Bare$. For $b:\Bare$, write the fiber $\Fib(b)=\{\,x\in\Real \mid \compare(x)=b\,\}$.
The \emph{residual kernel} collects exactly which distinct realizations $\compare$ identifies; injectivity means the kernel is empty; failure of injectivity means nontrivial kernel structure---not an informal ``remainder,'' but an object with fiberwise clique geometry, dual observables, refinement monotonicity, and functorial behavior under comparison morphisms.

\begin{center}
\fbox{\begin{minipage}{0.92\linewidth}
\textbf{Honesty (injective maps).} Not every comparison map is non-exhaustive or information-losing. Injective $\compare$ is a genuine countercase: the kernel is trivial and the sharp five-way equivalence collapses to the ``nothing to forget'' side. This paper proves the exact dichotomy: either the residual is trivial, or the full residual geometry of Theorem~\ref{thm:fund-res} (\S\ref{sec:fund-thm}) appears.
\end{minipage}}
\end{center}

\begin{center}
\fbox{\begin{minipage}{0.92\linewidth}
\textbf{Headline statements (preview)}\\[0.35em]
\textbf{(1) Fundamental equivalence.} Non-exhaustion, non-injectivity, witness diversity, fiber automorphism, and nontrivial residual kernel are equivalent; either all hold or none do (Theorem~\ref{thm:fund-equiv}).\\
\textbf{(2) Residual geometry.} The kernel decomposes fiberwise into complete indistinguishability cliques; observables are dual to kernel separation; refinement is graph coloring with exact per-fiber cost (Theorem~\ref{thm:fund-res}, Theorem~\ref{thm:res-complex}).\\
\textbf{(3) Universal forgetting.} The package is stated for arbitrary types and arbitrary $\compare$; certification language is an interpretation, not a hypothesis (\S\ref{sec:universal}).
\end{minipage}}
\end{center}

\medskip
\noindent
This paper develops that universal theory in full: what the map forgets, how it is organized, how much refinement costs, and when the structure is forced.

The motivating picture is certification ($\Bare$ bare certificates, $\Real$ realized witnesses), but the theorems need only $\compare$. Group homomorphisms, quotients, projections, measurements, compilers, and hashes inherit the same geometry because they are maps.

The results ascend through a sequence of summits:

\begin{enumerate}[label=\textbf{Summit \Roman*.},leftmargin=*]
\item \textbf{Non-exhaustion exists.} Adequate systems---those carrying
      witness-diversity data---force non-exhaustion as a theorem, not an
      assumption (\S\ref{sec:adequacy}).
\item \textbf{The residual has a canonical object.} The \emph{residual kernel}
      $\Ker(\compare)$ is the set of distinct realization pairs sharing a bare
      certificate. It is nonempty if and only if the system is non-exhaustive
      (\S\ref{sec:kernel}).
\item \textbf{The residual has geometry.} Fibers are the connected components of
      the kernel graph, and each fiber is a complete indistinguishability clique.
      Observables and refinements are governed by kernel structure
      (\S\S\ref{sec:fiber-geom}--\ref{sec:observable}).
\item \textbf{Resolution has exact complexity.} Refinement acts as graph coloring
      of fiber cliques. Resolving a fiber of size~$n$ requires exactly~$n$ colors.
      Terminal resolution is singletonization of refined classes
      (\S\ref{sec:refinement}).
\item \textbf{The fundamental equivalence.} Non-exhaustion, non-injectivity,
      witness diversity, fiber automorphism, and nontrivial kernel are five
      equivalent characterizations of the same phenomenon. Either all hold, or
      none do (\S\ref{sec:fund-equiv}).
\item \textbf{Universality.} The entire theory applies to any function between
      any two types. ``Certification'' is a motivating interpretation, not a
      formal hypothesis. The residual geometry is a universal theory of
      information loss (\S\ref{sec:universal}).
\item \textbf{Representation as structured forgetting (interpretive).}
      Section~\ref{sec:representation-as-structured-forgetting} traces how the
      fiber geometry of non-injective maps connects to selective articulation
      at the levels of function and ontology, and to reflexive non-exhaustion;
      \S\ref{sec:representation-conservation} adds the dual picture---conservation
      as class-level memory alongside map-level forgetting---with Lean interfaces
      in \emph{reflexive-closure-lean} (Remark~\ref{rem:structured-forgetting-lean}).
      Interpretive scope is fixed in Remark~\ref{rem:representation-interpretive}.
\end{enumerate}

\noindent
Expository order follows the ascent from existence to geometry to resolution;
the logical organizing theorem is the Classification Theorem in
\S\ref{sec:fund-equiv}, which governs all five summits simultaneously.

\noindent
All theorems are formalized in Lean~4~\cite{Lean4} with Mathlib~\cite{Mathlib} dependencies, building
cleanly with zero \texttt{sorry}. The formalization currently comprises approximately 235 theorems across 26 modules.

\subsection{Main theorem package}

The paper establishes the following five named results, each holding
unconditionally for every function between any two types. Logically, the
Classification Theorem governs the whole theory; expositionally, we ascend
through the canonical object, its geometry, and its resolution before
stating the equivalence.

\begin{enumerate}[label=(\Alph*),nosep]
\item \textbf{Classification Theorem}
      (Theorem~\ref{thm:fund-equiv}): non-exhaustion, non-injectivity, witness
      diversity, fiber automorphism, and nontrivial kernel are equivalent.
      Either all hold, or none do.
\item \textbf{Vanishing Criterion} (Theorem~\ref{thm:exhaust}): the system is
      exhaustive iff the kernel is empty iff every fiber is a subsingleton.
\item \textbf{Residual Geometry Theorem}
      (Theorem~\ref{thm:fund-res}): the residual kernel decomposes fiberwise
      into complete cliques; observables are dual to kernel separation;
      refinement resolves iff it separates every kernel pair; resolution
      decomposes into independent per-fiber coloring problems.
\item \textbf{Resolution Complexity Theorem}
      (Theorem~\ref{thm:res-complex}): resolving a fiber of size~$n$ requires
      exactly~$n$ distinct refinement values on that fiber.
\item \textbf{Certification--Refinement Gap Theorem}
      (Theorem~\ref{thm:cert-gap}): certification cannot resolve its own
      residual; refinement can, but at measurable cost; for adequate systems
      the cost is always positive.
\end{enumerate}

\noindent
These are not five independent results but a single coherent theory: (A)
classifies the phenomenon, (B) characterizes its absence, (C) describes its
geometry, (D) quantifies the cost of eliminating it, and (E) proves that the
cost is real and irreducible for adequate systems. Because the theory is
stated for an arbitrary function between arbitrary types, it applies to any
map that compresses structure (\S\ref{sec:universal}).

\noindent\textit{Logical vs.\ expository order.}
Logically, the Classification Theorem governs the whole theory; expositionally,
we ascend through the canonical object, its geometry, and its resolution
before returning to that equivalence as the governing law.

\begin{center}
\begin{tabular}{ll}
\toprule
\textbf{Question} & \textbf{Answer} \\
\midrule
What is information loss? & Nontrivial residual kernel \\
When does it occur? & Iff any of five equivalent conditions hold \\
Where does it live? & In the fibers (complete kernel cliques) \\
What can see it? & Exactly the non-bare-determined predicates \\
How is it resolved? & By refinements separating kernel pairs \\
What does resolution cost? & Exactly fiber-size-many colors \\
When is nothing left? & Exactly when $\compare$ is injective \\
Can the map see its own loss? & No (meta-level invisibility) \\
\bottomrule
\end{tabular}
\end{center}

\subsection{Novelty by synthesis}

The novelty of this paper is not the isolated appearance of fibers, kernels,
graph colorings, or factorization lemmas---each of those ingredients has
antecedents. The novelty is that these ingredients form \emph{one exact,
machine-checked theory}: a classification theorem for information loss, a
canonical residual object, a duality between observables and kernel separation,
an exact resolution criterion, a quantitative complexity law, and source
discharges from independent theorem programs---all valid for arbitrary maps,
all formalized with zero \texttt{sorry}. No previous work has assembled this
particular combination into a unified theory, and the assembly is not routine:
the Classification Theorem requires a classical construction (fiber swap) that
closes the circle between five superficially unrelated conditions; the
Residual Geometry Theorem packages unconditional structure that exists for
every function; and the formal verification is not decoration but the
mechanism by which the synthesis is checked. The central claim is therefore
not merely that non-injective maps have fibers, but that arbitrary maps carry
a full residual geometry with a classification theorem, a canonical residual
object, an observable duality, and an exact resolution complexity law.

\subsection{Relationship to the earlier summit}

Earlier work in Infinity Compression established the \emph{first summit} of
the program: canonical certification does not exhaust reflective realization
in a concrete flagship setting. A companion synthesis paper~\cite{Spivack2026-Strata} then showed that
this phenomenon sits inside a larger three-layer architecture linking NEMS,
APS, and Infinity Compression---with bridge theorems, the non-erasure
principle, and machine-checked coherence axioms proved definitionally.

The present paper is the \emph{higher summit forced by that earlier work}.
It extracts from those source discharges the map-theoretic core that survives
after domain-specific structure is stripped away. What appeared first as
irreducible residue above canonical certification in a concrete program is
shown here to be an instance of a universal geometry of information loss
under arbitrary maps. Summit~1 established that non-exhaustion is real and
unavoidable; Summit~2 identifies the exact mathematical object it produces,
proves its complete geometry, and shows that the phenomenon belongs to all
functions, not just to the flagship architecture.

In that sense this paper does not supersede the earlier summit. It completes
its abstraction. The old flagship now appears as a source discharge of the
universal theory (\S\ref{sec:discharges}).

%% ========================================================================
\section{Why this is not merely bookkeeping about kernels}\label{sec:not-bookkeeping}
%% ========================================================================

The residual attached to a comparison map is not an arbitrary set-theoretic difference quotient ``for flavor.'' It is the canonical locus where distinct realizations become indistinguishable at the bare level, and it is rich enough to support (i) a Boolean algebra of bare-determined observables dual to kernel separation, (ii) an exact notion of resolution cost via chromatic complexity on fibers, and (iii) functorial transport when comparisons change. In other words, ``what maps forget'' is not a slogan: it is a theorem-bearing program with quantitative structure and universal scope.

%% ========================================================================
\section{Reflective certification systems}\label{sec:rcs}
%% ========================================================================

We work in a universe-polymorphic setting. Fix types $\Bare$ and $\Real$ in
a common universe~$u$. For the universal theory, an RCS is simply a
convenient packaging of an arbitrary map together with two inert fields
retained for source-level compatibility (see Remark~\ref{rem:pi-only}).

\begin{definition}[Reflective certification system]\label{def:rcs}
A \emph{reflective certification system} $S$ on carriers $(\Bare, \Real)$
consists of:
\begin{itemize}[nosep]
\item a \emph{comparison map} $\compare : \Real \to \Bare$,
\item a \emph{canonical region} $\mathrm{Canonical} : \Bare \to \mathrm{Prop}$,
\item a \emph{soundness predicate} $\mathrm{Sound} : \mathrm{Prop}$.
\end{itemize}
\end{definition}

\begin{remark}[What drives the universal theory]\label{rem:pi-only}
Every theorem in \S\S\ref{sec:fund-equiv}--\ref{sec:fund-thm} depends only
on the comparison map $\compare : \Real \to \Bare$. The additional fields
$\mathrm{Canonical}$ and $\mathrm{Sound}$ are retained because they are part
of the originating formal interface and matter for source discharges
(\S\ref{sec:discharges}), but they play no role in the universal results.
This is what makes the theory broad: it applies to any function between any
two types, regardless of what mathematical structures the carriers carry.
\end{remark}

In the universal reading of this paper, one may equivalently regard an RCS
as an arbitrary map packaged in the originating interface vocabulary; the
name ``reflective certification system'' is inherited from the source program
and is not a restriction on the mathematics.
We write $\compare_S$ or simply $\compare$ for the comparison map of~$S$.

\begin{definition}[Exhaustion predicates]\label{def:exhaust}
Let $S$ be an RCS with comparison map $\compare$.
\begin{itemize}[nosep]
\item $S$ is \emph{minimally exhaustive} (or simply \emph{exhaustive}) if $\compare$
      is injective: $\compare(x) = \compare(y) \Rightarrow x = y$.
\item $S$ is \emph{strongly exhaustive} if $\compare$ admits a global section:
      $\exists\, s : \Bare \to \Real,\; \compare \circ s = \mathrm{id}$.
\item $S$ is \emph{non-exhaustive} if there exist $x, y \in \Real$ with
      $x \neq y$ and $\compare(x) = \compare(y)$.
\end{itemize}
\end{definition}

\begin{definition}[Fiber]\label{def:fiber}
The \emph{fiber} of $\compare$ over $b \in \Bare$ is
$\Fib(b) = \compare^{-1}(b) = \{\, x \in \Real \mid \compare(x) = b \,\}$.
\end{definition}

\begin{proposition}[Classical dichotomy]\label{prop:dichotomy}
For any RCS $S$, exactly one of the following holds: $S$ is exhaustive, or $S$
is non-exhaustive.
\end{proposition}

\begin{proof}
By classical logic applied to $\mathrm{Injective}(\compare)$.
\end{proof}

\begin{example}[Toy non-exhaustive system]\label{ex:toy}
Let $\Bare = \mathsf{Bool}$, $\Real = \mathsf{Bool} \times \mathsf{Bool}$,
and $\compare = \mathrm{fst}$. Then
$\Fib(\mathsf{true}) = \{(\mathsf{true}, \mathsf{false}),\,
(\mathsf{true}, \mathsf{true})\}$ has two elements, so the system is
non-exhaustive. The residual kernel contains the pair
$((\mathsf{true}, \mathsf{false}),\, (\mathsf{true}, \mathsf{true}))$.
The swap $\tau(a,b) = (a, \lnot b)$ is a non-identity $\compare$-preserving
endomorphism, illustrating the fiber automorphism of
Theorem~\ref{thm:fund-equiv}\ref{it:fa}. Resolving requires a refinement
with at least two values on each nontrivial fiber.
\end{example}

\begin{example}[Identity collapse]\label{ex:identity}
Let $\Bare = \Real = \mathsf{Bool}$ and $\compare = \mathrm{id}$. Every
fiber is a singleton, the kernel is empty, and the system is exhaustive.
No adequate or reflectively sufficient structure exists on this system.
\end{example}

\begin{remark}[Lean formalization]\label{rem:lean-rcs}
\leavevmode\par\begin{itemize}[nosep,leftmargin=*,before=\raggedright]
\item \texttt{ReflectiveCertificationSystem} in \texttt{Universal/ReflectiveCertificationSystem.lean}
\item \texttt{MinimalExhaustive}, \texttt{StrongExhaustive}, \texttt{NonExhaustive} in \texttt{Universal/\allowbreak{}ExhaustionDefinitions.lean}
\item \texttt{minimal\_exhaustive\_or\_\allowbreak{}nonExhaustive\_classical} in \texttt{Universal/Dichotomy.lean}
\item \texttt{Fiber}, \texttt{IsSection}, \texttt{LiftableWith} in \texttt{Universal/FiberBasics.lean} and \texttt{Universal/SectionsAndObstructions.lean}
\end{itemize}
\end{remark}

%% ========================================================================
\section{The fundamental equivalence theorem}\label{sec:fund-equiv}
%% ========================================================================

The deepest result of this paper is that five superficially different conditions
on an RCS are logically equivalent. We state it early because it governs
everything that follows.

\begin{theorem}[Classification Theorem]\label{thm:fund-equiv}
For any RCS $S = (\Bare, \Real, \compare)$, the following are equivalent:
\begin{enumerate}[label=(\roman*)]
\item\label{it:ne} \emph{Non-exhaustion.}
      $\exists\, x, y \in \Real,\; x \neq y \;\land\; \compare(x) = \compare(y)$.
\item\label{it:ni} \emph{Non-injectivity.}
      $\lnot\,\mathrm{Injective}(\compare)$.
\item\label{it:fa} \emph{Fiber automorphism.}
      $\exists\, \tau : \Real \to \Real,\;
      \compare \circ \tau = \compare \;\land\; \tau \neq \mathrm{id}$.
\item\label{it:wd} \emph{Witness diversity.}
      $\exists\, w_1, w_2 \in \Real,\; w_1 \neq w_2 \;\land\;
      \compare(w_1) = \compare(w_2)$.
\item\label{it:ker} \emph{Nontrivial kernel.}
      $\Ker(\compare) \neq \varnothing$.
\end{enumerate}
\end{theorem}

\begin{proof}
The equivalences \ref{it:ne}$\Leftrightarrow$\ref{it:ni},
\ref{it:ne}$\Leftrightarrow$\ref{it:wd}, and
\ref{it:ne}$\Leftrightarrow$\ref{it:ker} are definitional: non-exhaustion
\emph{is} witness diversity (the definitions coincide), non-exhaustion is the
negation of injectivity by the classical dichotomy, and the kernel is nonempty
exactly when a witness pair exists. These equivalences are recorded for
completeness; the theorem's content lies in the non-trivial direction.

The substantive direction is \ref{it:ne}$\Rightarrow$\ref{it:fa}. Given
$x \neq y$ with $\compare(x) = \compare(y)$, define the \emph{fiber swap}
\[
  \tau(z) \;=\; \begin{cases}
    y & \text{if } z = x, \\
    x & \text{if } z = y, \\
    z & \text{otherwise.}
  \end{cases}
\]
Then $\compare(\tau(z)) = \compare(z)$ for all~$z$ (checking the three cases
using $\compare(x) = \compare(y)$), and $\tau \neq \mathrm{id}$ because
$\tau(x) = y \neq x$. This construction uses classical decidability of equality
on $\Real$ (via \texttt{Classical.em} in the Lean formalization). The converse
direction \ref{it:fa}$\Rightarrow$\ref{it:ne} is constructive.

The converse \ref{it:fa}$\Rightarrow$\ref{it:ne}: if $\tau \neq \mathrm{id}$,
there exists $x$ with $\tau(x) \neq x$; then $\compare(\tau(x)) = \compare(x)$
with $\tau(x) \neq x$, witnessing non-exhaustion.
\end{proof}

The direction \ref{it:ne}$\Rightarrow$\ref{it:fa} is the \emph{fiber automorphism
diagonal}: the comparison geometry of any non-exhaustive system \emph{forces}
the existence of a non-trivial $\compare$-preserving self-map. The system's own
structure generates the symmetry that prevents exhaustion.

Theorem~\ref{thm:fund-equiv} is the \emph{Classification Theorem for
information loss under a map}. The equivalence of the five conditions is not
obvious \emph{a priori}: non-injectivity is a property of a function,
witness diversity is data carried by a structure, fiber automorphism is
an endomorphism condition, and kernel nonemptiness is a set-theoretic
statement. That these four perspectives collapse into a single
dichotomy---and that the fiber automorphism is \emph{forced} by
non-exhaustion via a classical construction---is the theorem's content.
Non-exhaustion is not merely a yes/no property; it is a condition with
five equivalent faces, each revealing a different aspect of the same
underlying geometry.

\begin{theorem}[All-or-nothing dichotomy]\label{thm:all-or-nothing}
For any RCS $S$, either all five conditions in
Theorem~\ref{thm:fund-equiv} hold simultaneously, or none do.
There is no intermediate position.
\end{theorem}

\begin{remark}[Honest ceiling]\label{rem:meta}
By the equivalence \ref{it:ne}$\Leftrightarrow$\ref{it:ni}, any predicate
$P$ satisfying $P(S) \Rightarrow \text{NonExhaustive}(S)$ also satisfies
$P(S) \Rightarrow \lnot\,\text{Injective}(\compare_S)$. This is logically
immediate, but its \emph{methodological} content is significant: at the level
of bare comparison maps, every forcing predicate must imply non-injectivity.
There is no way to force non-exhaustion without implying failure of
injectivity. So any summit class must carry structural data (witnesses,
twists, cardinality gaps) rather than a clever predicate on $\compare$ alone.
\end{remark}

\begin{remark}[Lean formalization]\label{rem:lean-equiv}
\leavevmode\par\begin{itemize}[nosep,leftmargin=*,before=\raggedright]
\item \texttt{equiv\_1\_iff\_2} through \texttt{equiv\_1\_iff\_5}, \texttt{all\_or\_nothing}, \texttt{FiveWayEquivalence}, \texttt{fiveWayEquivalence}, \texttt{nonExhaustive\_implies\_\allowbreak{}fiberAutomorphism} in \texttt{Summit/FundamentalEquivalence.lean}
\item \texttt{forcing\_implies\_not\_injective} in \texttt{Summit/Adequacy.lean}
\end{itemize}
\end{remark}

%% ========================================================================
\section{The exhaustion criterion}\label{sec:exhaust-crit}
%% ========================================================================

The exhaustion criterion is the exact dual vanishing theorem to the
fundamental equivalence: where Theorem~\ref{thm:fund-equiv} classifies
the \emph{presence} of certification failure, the exhaustion criterion
classifies its \emph{absence}.

\begin{theorem}[Vanishing Criterion]\label{thm:exhaust}
For any RCS $S$, the following are equivalent:
\begin{enumerate}[label=(\roman*)]
\item $\compare$ is injective.
\item $\Ker(\compare) = \varnothing$.
\item Every fiber $\Fib(b)$ is a subsingleton.
\item $S$ is not non-exhaustive.
\end{enumerate}
\end{theorem}

There is no partial exhaustion: the system is completely exhaustive if and only
if the kernel is empty, if and only if every fiber has at most one element.

\begin{remark}[Lean formalization]
\leavevmode\par\begin{itemize}[nosep,leftmargin=*,before=\raggedright]
\item \texttt{exhaustive\_iff\_kernel\_empty}, \texttt{kernel\_empty\_iff\_\allowbreak{}all\_fibers\_subsingleton}, \texttt{exhaustive\_iff\_not\_nonExhaustive} in \texttt{Residual/FundamentalTheorem.lean}
\end{itemize}
\end{remark}

%% ========================================================================
\section{The adequacy summit}\label{sec:adequacy}
%% ========================================================================

The fundamental equivalence tells us \emph{what} non-exhaustion is equivalent to.
The adequacy summit tells us \emph{when} it is forced.

\begin{definition}[Adequate reflective system]\label{def:adequate}
An \emph{adequate reflective system} is an RCS equipped with
\emph{witness diversity data}: two distinguished elements
$w_1, w_2 \in \Real$ with $w_1 \neq w_2$ and
$\compare(w_1) = \compare(w_2)$.
\end{definition}

This is \emph{data}, not a proposition: the witnesses are fields of the
structure, not existentially quantified. The distinction matters because it
gives the class computational content.

\begin{theorem}[Adequacy Theorem]\label{thm:adequate-ne}
Every adequate reflective system is non-exhaustive.
\end{theorem}

\begin{proof}
The two witnesses are the existential witnesses for non-exhaustion.
\end{proof}

\begin{proposition}[Non-circularity]\label{prop:non-circ}
The axioms of adequacy mention neither non-exhaustion nor injectivity.
Identity-style collapse ($\compare = \mathrm{id}$) provably cannot satisfy
the axioms: $\compare(w_1) = \compare(w_2)$ with $\compare = \mathrm{id}$
forces $w_1 = w_2$, contradicting distinctness. More generally, any adequate
system has non-injective comparison.
\end{proposition}

\begin{definition}[Reflectively sufficient system]\label{def:refl-suff}
A \emph{reflectively sufficient system} is an RCS equipped with a
non-identity $\compare$-preserving endomorphism
$\tau : \Real \to \Real$ satisfying $\compare \circ \tau = \compare$
and $\tau \neq \mathrm{id}$.
\end{definition}

\begin{theorem}[Self-generation]\label{thm:self-gen}
Every reflectively sufficient system is adequate. Since $\tau \neq \mathrm{id}$,
there exists $x$ with $\tau(x) \neq x$; then $(x, \tau(x))$ is a witness pair
with $\compare(x) = \compare(\tau(x))$.
\end{theorem}

The self-generation theorem shows that witness diversity need not be externally
supplied: the system's own internal symmetry generates it. This is the
\emph{Gödelian core} of the program. In Gödel's incompleteness theorem~\cite{Godel1931}, a
sufficiently rich formal system's own expressive power constructs a sentence
the system cannot decide. Here, a non-exhaustive system's own comparison
geometry constructs the endomorphism that witnesses its own non-exhaustion.
The mechanism is the same: self-action produces the structure that prevents
completeness.

By the fundamental equivalence (Theorem~\ref{thm:fund-equiv}), the converse
also holds: every non-exhaustive system admits a $\compare$-preserving
endomorphism (via the fiber swap construction). So the three
classes---non-exhaustive, adequate, reflectively sufficient---coincide
propositionally. The complete self-generation chain is:
\[
  \text{ReflectivelySufficient}
  \;\xrightarrow{\text{twist produces witnesses}}\;
  \text{Adequate}
\]
\[
  \text{Adequate}
  \;\xrightarrow{\text{witnesses force non-exhaustion}}\;
  \text{NonExhaustive}
\]
\[
  \text{NonExhaustive}
  \;\xrightarrow{\text{fiber swap}}\;
  \text{ReflectivelySufficient}.
\]
The circle closes: there is no escape from any of the three without escaping
all of them.

\begin{remark}[Lean formalization]
\leavevmode\par\begin{itemize}[nosep,leftmargin=*,before=\raggedright]
\item \texttt{AdequateReflectiveSystem}, \texttt{adequate\_nonExhaustive}, \texttt{adequate\_not\_injective}, \texttt{identityCompare\_not\_adequate} in \texttt{Summit/Adequacy.lean}
\item \texttt{ReflectivelySufficientSystem}, \texttt{toAdequate}, \texttt{reflectivelySufficient\_\allowbreak{}nonExhaustive} in \texttt{Summit/SelfGeneration.lean}
\item IC discharge: \texttt{icAdequateReflectiveSystem} in \texttt{Summit/ICAdequateDischarge.lean}
\end{itemize}
\end{remark}

\begin{remark}[Certification versus refinement]\label{rem:cert-vs-refine}
A potential confusion must be addressed. The resolution theory
(\S\ref{sec:refinement}) shows that the residual can always be eliminated
by adding a sufficiently fine refinement. Does this contradict the claim
that certification does not exhaust realization?

No. The distinction is between the \emph{certification itself} (the fixed
comparison map $\compare$) and \emph{supplementary data} (a refinement
$r : \Real \to E$ added on top of $\compare$). For adequate systems,
$\compare$ is provably non-injective---the residual is intrinsic to the
certification and cannot be eliminated by the certification itself
(Proposition~\ref{prop:non-circ}). A refinement can resolve the residual,
but the refinement is extra structure \emph{beyond} the certification.
The identity refinement $r = \mathrm{id}$ always resolves, but it means
``remember all of $\Real$''---defeating the purpose of certification,
which is to compress realization into something smaller.

The precise statement is the \emph{certification-refinement gap theorem}:
for any adequate system, (i)~$\compare$ is not injective, (ii)~there
exists a refinement that resolves, and (iii)~no subsingleton-valued
refinement suffices. The residual is real, its resolution requires
genuine extra data, and the cost is measurable.
\end{remark}

%% ========================================================================
\section{The residual kernel}\label{sec:kernel}
%% ========================================================================

\begin{definition}[Residual kernel]\label{def:kernel}
The \emph{residual kernel} of an RCS $S$ is
\[
  \Ker(\compare) \;=\; \bigl\{\,(x,y) \in \Real \times \Real
    \;\big|\; x \neq y \;\land\; \compare(x) = \compare(y)\,\bigr\}.
\]
\end{definition}

The kernel records every ordered pair of distinct realizations that
certification cannot distinguish. It is a symmetric relation
($(x,y) \in \Ker \Leftrightarrow (y,x) \in \Ker$) and is intrinsic to
$\compare$: two systems with the same comparison map have identical kernels.

From this point onward, non-exhaustion should no longer be thought of merely
as the existence of two colliding realizations. Its canonical carrier is the
residual kernel, and the rest of the theory consists in reading geometry,
observability, and resolution from that single object.

\begin{theorem}[Kernel characterization]\label{thm:kernel-char}
\leavevmode
\begin{enumerate}[label=(\roman*),nosep]
\item $\Ker(\compare) \neq \varnothing$ if and only if $S$ is non-exhaustive.
\item $\Ker(\compare)$ depends only on $\compare$.
\item Two systems with the same $\compare$ have identical kernels, identical
      unresolved kernels for any refinement, and identical incompressibility
      at any type.
\item $\Ker(\compare)$ is the off-diagonal of the fiber product:
      $\Ker(\compare) = (\Real \times_\Bare \Real) \setminus \Delta$,
      where $\Real \times_\Bare \Real = \{(x,y) \mid \compare(x) = \compare(y)\}$
      and $\Delta = \{(x,x)\}$. The system is exhaustive iff the fiber product
      equals the diagonal.
\end{enumerate}
\end{theorem}

\begin{definition}[Veil]\label{def:veil}
For an adequate system $A$, the \emph{veil} is the residual kernel
$\Ker(\compare_A)$. By Theorem~\ref{thm:adequate-ne}, the veil is nonempty.
The adequate witnesses $(w_1, w_2)$ are a canonical element of the veil.
\end{definition}

\begin{theorem}[Veil persistence]\label{thm:veil}
For any adequate system, the trivial refinement (to a point) leaves the veil
fully intact: the unresolved kernel equals the full kernel.
\end{theorem}

\begin{remark}[Lean formalization]
\leavevmode\par\begin{itemize}[nosep,leftmargin=*,before=\raggedright]
\item \texttt{ResidualKernel}, \texttt{residualKernel\_nonempty\_iff}, \texttt{residualKernel\_symmetric}, \texttt{residualKernel\_eq\_of\_\allowbreak{}compare\_eq}, \texttt{residualTransfer\_\allowbreak{}unresolvedKernel} in \texttt{Residual/\allowbreak{}ResidualKernel.lean}
\item \texttt{adequate\_kernel\_nonempty}, \texttt{adequate\_witnesses\_in\_kernel}, \texttt{adequate\_veil\_persists\_trivial} in \texttt{Residual/ResidualKernel.lean}
\end{itemize}
\end{remark}

%% ========================================================================
\section{Fiber geometry}\label{sec:fiber-geom}
%% ========================================================================

The kernel decomposes fiberwise, and each fiber carries a maximally dense
internal structure.

\begin{definition}[Per-fiber kernel]\label{def:kernel-at}
For $b \in \Bare$, the \emph{per-fiber kernel} is
$\KerAt_b = \{(x,y) \mid x \neq y,\, \compare(x) = b,\, \compare(y) = b\}$.
This is the off-diagonal of $\Fib(b)$.
\end{definition}

\begin{theorem}[Fiberwise decomposition]\label{thm:fiber-decomp}
$\Ker(\compare) = \bigcup_{b \in \Bare} \KerAt_b$.
\end{theorem}

\begin{theorem}[Complete-fiber theorem]\label{thm:complete-fiber}
For every $b \in \Bare$, the per-fiber kernel $\KerAt_b$ is a
\emph{complete graph}: every pair of distinct elements in $\Fib(b)$
is a kernel edge.
\end{theorem}

\begin{proof}
If $x, y \in \Fib(b)$ with $x \neq y$, then $\compare(x) = b = \compare(y)$,
so $(x,y) \in \KerAt_b$.
\end{proof}

Fibers are not merely clusters of collisions---they are \emph{maximally dense
indistinguishability cliques}. Every possible collision within a fiber actually
occurs. There are no gaps in the residual within a fiber.

\begin{theorem}[Connectivity theorem]\label{thm:connectivity}
Define $x \sim y$ iff $\compare(x) = \compare(y)$ (fiber equivalence). Then:
\begin{enumerate}[label=(\roman*),nosep]
\item $\sim$ is an equivalence relation.
\item $x \sim y$ iff $x = y$ or $(x,y) \in \Ker(\compare)$.
\item The equivalence classes of $\sim$ are exactly the fibers $\Fib(b)$.
\item Fibers are the connected components of the kernel graph.
\end{enumerate}
\end{theorem}

\begin{proof}
Part~(i) is immediate. For~(ii), $x \sim y$ means $\compare(x) = \compare(y)$;
if $x \neq y$ this is a kernel pair, otherwise $x = y$. Part~(iii) follows from
the definition of fibers. Part~(iv) follows from~(ii)--(iii): two elements are
reachable via kernel edges iff they share a fiber.
\end{proof}

The kernel neighbors of any element $x$ are exactly
$\Fib(\compare(x)) \setminus \{x\}$---all other members of its fiber.

\begin{remark}[Lean formalization]
\leavevmode\par\begin{itemize}[nosep,leftmargin=*,before=\raggedright]
\item \texttt{KernelAt}, \texttt{residualKernel\_eq\_iUnion\_kernelAt}, \texttt{kernelAt\_complete}, \texttt{kernelAt\_eq\_offDiagonal\_fiber} in \texttt{Residual/KernelStratification.lean} and \texttt{Residual/MinimalResolution.lean}
\item \texttt{KernelRel}, \texttt{KernelRelRefl}, \texttt{kernelRelRefl\_equivalence}, \texttt{fiberEquiv\_iff\_eq\_or\_kernelRel}, \texttt{fiber\_eq\_kernelRelRefl\_class}, \texttt{kernelNeighbors\_eq\_fiber\_minus\_self} in \texttt{Residual/KernelGraph.lean} and \texttt{Residual/MinimalResolution.lean}
\end{itemize}
\end{remark}

%% ========================================================================
\section{Observable duality}\label{sec:observable}
%% ========================================================================

The kernel determines a sharp boundary between what certification can and
cannot see.

\begin{definition}[Bare-determined predicate]\label{def:bare-det}
A predicate $P : \Real \to \mathrm{Prop}$ is \emph{bare-determined} if it
factors through $\compare$: there exists $Q : \Bare \to \mathrm{Prop}$ with
$P(x) \Leftrightarrow Q(\compare(x))$ for all $x$.
\end{definition}

\begin{theorem}[Observable algebra]\label{thm:obs-algebra}
The bare-determined predicates form a Boolean subalgebra of the full predicate
algebra on $\Real$: they are closed under $\land$, $\lor$, $\lnot$, $\to$,
$\leftrightarrow$, and contain $\top$ and $\bot$.
\end{theorem}

\begin{proof}
If $P$ factors through $Q_P$ and $P'$ factors through $Q_{P'}$, then
$P \land P'$ factors through $Q_P \land Q_{P'}$, and similarly for the
other connectives.
\end{proof}

\begin{theorem}[Universal factorization]\label{thm:univ-factor}
A function $f : \Real \to \alpha$ (with $\alpha$ nonempty) is constant on
every fiber of $\compare$ if and only if it factors through $\compare$:
there exists $g : \Bare \to \alpha$ with $f = g \circ \compare$.
\end{theorem}

\begin{proof}
The forward direction constructs $g$ classically: for each $b \in \Bare$,
if $b$ is in the image of $\compare$, set $g(b) = f(x)$ for any $x$ with
$\compare(x) = b$ (well-defined by the fiber-constancy hypothesis);
otherwise, set $g(b)$ to an arbitrary element of $\alpha$.
The converse is immediate.
\end{proof}

This theorem says $\compare$ is the \emph{universal coarse observable}:
every observable that cannot see inside fibers must factor through $\compare$,
and anything coarser than $\compare$ is blind to all kernel pairs.

\begin{theorem}[Kernel-observable duality]\label{thm:duality}
A predicate $P : \Real \to \mathrm{Prop}$ is not bare-determined
if and only if it separates at least one kernel pair: there exist
$(x,y) \in \Ker(\compare)$ with $(P(x) \land \lnot P(y))$ or
$(\lnot P(x) \land P(y))$.
\end{theorem}

\begin{proof}
Forward: if $P$ is not bare-determined, it is not constant on some fiber
(otherwise it would factor through $\compare$ by Theorem~\ref{thm:univ-factor}
applied to the characteristic function). So there exist $x, y$ in the same
fiber with $P(x)$ and $\lnot P(y)$; since they share a fiber and are
distinguished by $P$, they are distinct, hence a kernel pair.

Backward: if $P$ separates $(x,y) \in \Ker$, then $P$ is not constant on
$\Fib(\compare(x))$, so it cannot factor through $\compare$.
\end{proof}

This theorem completely characterizes the boundary between
map-visible and map-invisible predicates. The residual
observables are not an indeterminate remainder; they are exactly the
predicates that cut across kernel pairs. The kernel and the non-bare
observables are dual objects: each determines the other.

\begin{remark}[Lean formalization]
\leavevmode\par\begin{itemize}[nosep,leftmargin=*,before=\raggedright]
\item \texttt{BareDetermined}, \texttt{bareDetermined\_and/or/not/imp}, \texttt{bareDetermined\_of\_comp} in \texttt{Residual/ObservableAlgebra.lean}
\item \texttt{constant\_on\_fibers\_iff\_factors}, \texttt{factors\_of\_constant\_on\_fibers} in the same file
\item \texttt{not\_bareDetermined\_iff\_separates\_kernel} in \texttt{Residual/ResidualKernel.lean}
\end{itemize}
\end{remark}

%% ========================================================================
\section{Refinement and resolution}\label{sec:refinement}
%% ========================================================================

A \emph{refinement} $r : \Real \to E$ provides extra information beyond
$\compare$. The \emph{refined comparison} is the pair
$(\compare, r) : \Real \to \Bare \times E$.

\begin{definition}[Resolution]\label{def:resolution}
A refinement $r$ \emph{resolves} the system if the refined comparison
$(\compare, r)$ is injective.
\end{definition}

\begin{theorem}[Resolution criterion]\label{thm:resolution}
A refinement resolves the system if and only if it separates every kernel pair:
for all $(x,y) \in \Ker(\compare)$, $r(x) \neq r(y)$.
\end{theorem}

\begin{proof}
If $(\compare, r)$ is injective and $(x,y) \in \Ker$, then
$\compare(x) = \compare(y)$, so injectivity forces $r(x) \neq r(y)$
(else $(\compare(x), r(x)) = (\compare(y), r(y))$ would give $x = y$,
contradicting $x \neq y$). Conversely, if $r$ separates every kernel pair
and $(\compare(x), r(x)) = (\compare(y), r(y))$, then either $x = y$
(done) or $(x,y) \in \Ker$ and $r(x) = r(y)$, contradicting separation.
\end{proof}

\begin{theorem}[Fiber-local resolution]\label{thm:local-res}
Resolution decomposes into independent per-fiber problems: a refinement
resolves the whole system if and only if it resolves every fiber individually.
Cross-fiber pairs are never kernel pairs.
\end{theorem}

\begin{theorem}[Coloring interpretation]\label{thm:coloring}
A refinement acts as a graph coloring of each fiber clique: it assigns
a color $r(x) \in E$ to each vertex. It resolves a fiber if and only
if it properly colors the complete graph on that fiber.
\end{theorem}

Since each fiber kernel is a complete graph $K_n$ on $n = |\Fib(b)|$ vertices,
the chromatic number is~$n$.

\begin{theorem}[Resolution complexity]\label{thm:res-complex}
Resolving a fiber of size $n$ requires at least $n$ distinct refinement values
on that fiber. Equivalently, the refinement must be injective on the fiber.
This is the chromatic number of the complete graph $K_n$: the exact cost of
eliminating the residual at a given fiber is its size.
\end{theorem}

\begin{definition}[Refined fiber class]\label{def:refined-class}
Given a refinement $r$, the \emph{refined fiber class} at $(b, e)$ is
$\{x \in \Real \mid \compare(x) = b \land r(x) = e\}$. The fiber $\Fib(b)$
is the disjoint union of its refined classes.
\end{definition}

\begin{theorem}[Terminal refinement]\label{thm:terminal}
A refinement is terminal (fully resolving) if and only if every refined class
is a singleton.
\end{theorem}

\begin{definition}[Refinement preorder]\label{def:refines-to}
Refinement $r_2$ \emph{refines} $r_1$ if equal $r_2$-values imply equal
$r_1$-values. This is a preorder with the trivial refinement (to a point)
as bottom and the identity refinement (to $\Real$) as top.
\end{definition}

\begin{definition}[Unresolved kernel]\label{def:unresolved}
The \emph{unresolved kernel} of a refinement $r$ is
$\Ker^{\mathrm{unres}}_r = \{(x,y) \in \Ker(\compare) \mid r(x) = r(y)\}$.
\end{definition}

\begin{theorem}[Monotonicity and progress]\label{thm:monotone}
\leavevmode
\begin{enumerate}[label=(\roman*),nosep]
\item The unresolved kernel is antitone: finer refinements have weakly smaller
      unresolved kernels.
\item The resolved kernel (complement in $\Ker$) is monotone: finer refinements
      resolve weakly more pairs.
\item A refinement resolves iff its unresolved kernel is empty.
\item Any refinement separating at least one new kernel pair strictly shrinks the
      unresolved kernel.
\end{enumerate}
\end{theorem}

\begin{theorem}[Trivial-resolves-iff-exhaustive]\label{thm:triv-res}
The trivial refinement resolves if and only if the system is exhaustive.
The cheapest possible resolution works exactly when there is no residual.
\end{theorem}

\begin{theorem}[Adequate resolution bound]\label{thm:adequate-bound}
For adequate systems: subsingleton-valued refinements cannot resolve.
The identity refinement always resolves. So the resolution complexity of an
adequate system is bounded below by~$2$ and above by~$|\Real|$.
\end{theorem}

\begin{definition}[Nontrivial locus]\label{def:nontrivial-locus}
The \emph{nontrivial locus} of an RCS is the set of bare certificates whose
fibers contain at least two distinct elements. The system is exhaustive iff
the nontrivial locus is empty; non-exhaustive iff it is nonempty. The
nontrivial locus is nonempty iff the kernel is nonempty.
\end{definition}

The nontrivial locus is the ``support'' of the residual: it tells you
\emph{which} certificates carry hidden structure. The resolution cost at
each certificate in the locus is the fiber size at that certificate.

\begin{definition}[Persistent observable]\label{def:persistent}
A non-bare-determined predicate is \emph{persistent} under a refinement $r$
if it separates a kernel pair that $r$ also fails to separate (i.e., a pair
in the unresolved kernel of $r$).
\end{definition}

\begin{theorem}[Persistence dichotomy]\label{thm:persist}
Either a refinement resolves (and no observable is persistent) or the
unresolved kernel is nonempty (and persistent observables exist).
\end{theorem}

Persistent observables are the predicates that remain ``invisible to
certification'' even after the refinement is applied. They measure what
the refinement has not yet resolved.

\begin{theorem}[Certification-refinement gap]\label{thm:cert-gap}
For any adequate system $A$:
\begin{enumerate}[label=(\roman*),nosep]
\item $\compare_A$ is not injective (the certification has a residual).
\item There exists a refinement that resolves (the identity refinement).
\item No subsingleton-valued refinement resolves (the resolution cost is $\geq 2$).
\end{enumerate}
The residual is intrinsic to the certification level: it is determined by
$\compare$ alone and cannot be changed by altering other system data. Eliminating
it requires genuine extra data beyond the certification.
\end{theorem}

\begin{remark}[Lean formalization]
\leavevmode\par\begin{itemize}[nosep,leftmargin=*,before=\raggedright]
\item Key definitions: \texttt{Refinement}, \texttt{refinedCompare}, \texttt{RefinesTo}, \texttt{UnresolvedKernel}, \texttt{RefinedFiberClass}
\item Key theorems: \texttt{resolves\_iff\_separates\_\allowbreak{}all\_kernel\_pairs}, \texttt{resolves\_iff\_properly\_\allowbreak{}colors\_all}, \texttt{terminal\_iff\_all\_\allowbreak{}singletons}, \texttt{trivial\_resolves\_iff\_exhaustive}
\item Spread across \texttt{Residual/ResidualStructure.lean}, \texttt{Residual/MinimalResolution.lean}, \texttt{Residual/ResolutionComplexity.lean}, \texttt{Residual/QuantitativeResidual.lean}, \texttt{Residual/RefinementOrder.lean}, \texttt{Residual/Incompressibility.lean}, \texttt{Residual/KernelStratification.lean}
\item Fiber spectrum: \texttt{FiberTrivial}, \texttt{FiberNontrivial}, \texttt{NontrivialLocus} in \texttt{Residual/FiberSpectrum.lean}
\item Categorical kernel: \texttt{FiberProduct}, \texttt{Diagonal}, \texttt{residualKernel\_eq\_\allowbreak{}fiberProduct\_minus\_diagonal} in \texttt{Residual/CategoricalKernel.lean}
\item Appendix~\ref{app:index} records the main paper-to-Lean correspondence for this development.
\end{itemize}
\end{remark}

%% ========================================================================
\section{Optimal certification and residual measurement}\label{sec:optimal}
%% ========================================================================

The resolution theory raises a natural question: if refinement can always
eliminate the residual, does anything genuinely remain? The answer depends
on what ``remains'' means.

\begin{theorem}[Optimal refinement existence]\label{thm:opt-exist}
For any RCS $S$, there exists a refinement that fully resolves the system:
the identity refinement $r = \mathrm{id} : \Real \to \Real$ makes the
refined comparison $(\compare, \mathrm{id})$ injective.
\end{theorem}

\begin{proof}
If $(\compare(x), x) = (\compare(y), y)$, then $x = y$.
\end{proof}

So the residual can always be eliminated. But the identity refinement means
``remember all of $\Real$''---it defeats the purpose of certification, which
is to compress realization into something smaller. The real question is the
\emph{cost} of resolution.

\begin{theorem}[Zero residual characterization]\label{thm:zero-res}
For any RCS $S$, the following are equivalent:
\begin{enumerate}[label=(\roman*),nosep]
\item $\Ker(\compare) = \varnothing$ (kernel empty).
\item $\compare$ is injective.
\item Every fiber is trivial (at most one element).
\item The nontrivial locus is empty.
\item The trivial refinement (to a point) resolves.
\item Every predicate on $\Real$ is bare-determined.
\end{enumerate}
When any of these hold, the residual is completely absent: certification sees
everything, no extra data is needed, and there is no hidden structure.
\end{theorem}

\begin{definition}[Refinement gap]\label{def:gap}
The \emph{refinement gap} of a refinement $r$ is its unresolved kernel:
$\mathrm{Gap}(r) = \{(x,y) \in \Ker(\compare) \mid r(x) = r(y)\}$.
This measures how far $r$ is from full resolution.
\end{definition}

\begin{theorem}[Gap properties]\label{thm:gap-props}
\leavevmode
\begin{enumerate}[label=(\roman*),nosep]
\item The gap is empty iff the refinement resolves.
\item The gap of the trivial refinement is the full kernel.
\item The gap of the identity refinement is empty.
\item Finer refinements have smaller gaps (monotonicity).
\end{enumerate}
\end{theorem}

The gap gives a precise ``distance to resolution'': it is the set of kernel
pairs that the refinement has not yet separated. Each step of refinement
that separates a new kernel pair strictly shrinks the gap
(Theorem~\ref{thm:monotone}(iv)).

\begin{theorem}[Adequate irreducible residual]\label{thm:adequate-irred}
For any adequate system $A$:
\begin{enumerate}[label=(\roman*),nosep]
\item The gap of the trivial refinement is nonempty (it equals the full
      kernel, which is nonempty by adequacy).
\item The gap of any subsingleton-valued refinement is nonempty.
\item The adequate witnesses $(w_1, w_2)$ are in the gap of every
      subsingleton-valued refinement.
\end{enumerate}
The residual of an adequate system has irreducible depth: it cannot be
eliminated without genuine extra data.
\end{theorem}

\begin{theorem}[Nothing remains at optimal]\label{thm:nothing-remains}
If $\compare$ is injective, then:
\begin{enumerate}[label=(\roman*),nosep]
\item $\Ker(\compare) = \varnothing$.
\item The nontrivial locus is empty.
\item Every predicate on $\Real$ is bare-determined---certification sees
      everything.
\item The trivial refinement resolves---no extra data is needed.
\end{enumerate}
At optimal certification, the veil is gone. There is no hidden structure.
\end{theorem}

The picture is now complete. Certification determines a comparison map.
If the map is injective, the residual is absent and certification exhausts
realization. If it is not---as the adequacy theorem forces for systems with
witness diversity---the residual is a nontrivial geometric object: the kernel,
decomposed into complete fiber cliques, with exact resolution cost. The
residual can be resolved by adding extra data, but the cost is real and
measurable, and for adequate systems it is always positive.

\begin{remark}[Lean formalization]
\leavevmode\par\begin{itemize}[nosep,leftmargin=*,before=\raggedright]
\item \texttt{optimal\_refinement\_exists}, \texttt{zero\_residual\_iff\_injective}, \texttt{zero\_residual\_iff\_all\_trivial}, \texttt{zero\_residual\_iff\_trivial\_resolves} in \texttt{Residual/OptimalCertification.lean}
\item \texttt{refinementGap}, \texttt{gap\_empty\_iff\_resolves}, \texttt{gap\_trivial\_eq\_kernel}, \texttt{gap\_identity\_empty}, \texttt{gap\_antitone} in the same file
\item \texttt{adequate\_irreducible\_residual}, \texttt{nothing\_remains\_at\_\allowbreak{}optimal}, \texttt{certification\_refinement\_gap} in the same file
\end{itemize}
\end{remark}

%% ========================================================================
\section{Functoriality and invariance}\label{sec:functorial}
%% ========================================================================

The residual geometry is intrinsic to the comparison map and transforms
coherently under maps between systems.

\begin{definition}[Comparison-compatible map]\label{def:compat}
A \emph{comparison-compatible map} from $S_1$ to $S_2$ (on the same carriers)
is a function $\varphi : \Real \to \Real$ with
$\compare_2 \circ \varphi = \compare_1$.
\end{definition}

\begin{theorem}[Functoriality]\label{thm:functorial}
\leavevmode
\begin{enumerate}[label=(\roman*),nosep]
\item An injective compatible map sends kernel pairs to kernel pairs and
      preserves fiber membership.
\item Compatible maps compose: if $\varphi : S_1 \to S_2$ and
      $\psi : S_2 \to S_3$, then $\psi \circ \varphi : S_1 \to S_3$.
\item The identity is compatible with itself.
\item Composition is associative, and the identity is a two-sided unit.
\item Two systems with the same $\compare$ have identical kernels,
      identical unresolved kernels, and identical incompressibility.
\end{enumerate}
Thus RCS on fixed carriers and compatible maps form a category.
\end{theorem}

\begin{theorem}[Optimal coarsening]\label{thm:optimal}
$\compare$ is the finest fiber-constant map: any function constant on fibers
factors through $\compare$. Any map coarser than $\compare$ is blind to all
kernel pairs.
\end{theorem}

\begin{theorem}[Witness independence]\label{thm:witness-indep}
Fiber nontriviality, the fiber partition, and the bare-determined algebra all
depend only on $\compare$, not on the choice of witnesses, canonical region,
or soundness predicate.
\end{theorem}

\begin{remark}[Lean formalization]
\leavevmode\par\begin{itemize}[nosep,leftmargin=*,before=\raggedright]
\item \texttt{CompareCompatible}, \texttt{compCompatible}, \texttt{idCompatible}, \texttt{compCompatible\_assoc}, \texttt{compatible\_preserves\_kernel}, \texttt{CoarserThanCompare}, \texttt{compare\_is\_finest\_fiber\_constant}, \texttt{coarser\_blind\_to\_kernel} in \texttt{Residual/FundamentalTheorem.lean}
\item \texttt{nonExhaustive\_depends\_only\_on\_compare}, \texttt{fiber\_eq\_of\_compare\_eq}, \texttt{bareDetermined\_depends\_only\_on\_compare} in \texttt{Residual/WitnessFiber.lean}
\end{itemize}
\end{remark}

%% ========================================================================
\section{The fundamental theorem of residual geometry}\label{sec:fund-thm}
%% ========================================================================

We package the entire residual theory into a single statement.

\begin{theorem}[Residual Geometry Theorem]\label{thm:fund-res}
For any comparison map $\compare : \Real \to \Bare$, the results of
\S\S\ref{sec:kernel}--\ref{sec:refinement} hold unconditionally: the kernel
characterization (Theorem~\ref{thm:kernel-char}), fiberwise decomposition
(Theorem~\ref{thm:fiber-decomp}), complete-fiber structure
(Theorem~\ref{thm:complete-fiber}), kernel-observable duality
(Theorem~\ref{thm:duality}), resolution criterion
(Theorem~\ref{thm:resolution}), fiber-local resolution
(Theorem~\ref{thm:local-res}), resolution complexity
(Theorem~\ref{thm:res-complex}), and the exhaustion criterion
(Theorem~\ref{thm:exhaust}). No hypotheses on the RCS are required.
\end{theorem}

The adequacy and self-generation theorems (\S\ref{sec:adequacy}) tell us
\emph{when} the kernel is nontrivial; the fundamental theorem tells us
\emph{what} the kernel is and how it behaves, for every system.

\begin{remark}[Lean formalization]
\leavevmode\par\begin{itemize}[nosep,leftmargin=*,before=\raggedright]
\item \texttt{FundamentalResidualPackage}, \texttt{fundamentalResidualPackage} in \texttt{Residual/FundamentalTheorem.lean}
\item The package is a Lean structure whose fields are the eight component theorems; the constructor \texttt{fundamentalResidualPackage} builds it for any RCS
\end{itemize}
\end{remark}

%% ========================================================================
\section{Source discharges}\label{sec:discharges}
%% ========================================================================

Historically, the universal theory did not arise by beginning with arbitrary
maps and asking abstract questions about fibers. It was forced by concrete
theorems in three source programs. Infinity Compression exhibited irreducible
residue above canonical certification; NEMS supplied a computability-facing
non-erasure barrier; APS isolated a local-to-global realization architecture
that prevented the synthesis from collapsing into a vacuous slogan. The
present paper extracts from those discharges the universal map-theoretic core
that survives after all domain-specific structure is stripped away.

Historically, the theory arose from these source programs; formally, they
now serve as concrete discharges of a universal map-theoretic framework.
The abstract theory developed in \S\S\ref{sec:fund-equiv}--\ref{sec:fund-thm}
is universal: it applies to any comparison map. But the theory acquires force
when discharged against concrete mathematical architectures that exhibit
non-exhaustion for principled reasons. The program originated in a synthesis
of three independent research lines, each governing one layer of a reflexive
architecture.

\paragraph{The three-layer origin.}
The Strata program unifies three lines into a single architecture:

\begin{center}
\begin{tabular}{lll}
\toprule
Layer & Source program & Structural role \\
\midrule
Outer & NEMS & Admissibility of determinacy under closure \\
Middle & APS & Exact mechanics of uniform realization \\
Inner & Infinity Compression & Reflective residue above canonical certification \\
\bottomrule
\end{tabular}
\end{center}

\noindent
The \emph{Non-Erasure Principle} proved in the Strata formalization states
that under explicit cross-layer coherence hypotheses, outer semantic
non-exhaustion and inner enriched irreducibility hold simultaneously and
resist joint flattening. The \emph{Unified Non-Erasure Law} gives the
biconditional: enriched irreducibility $\Leftrightarrow$ existence of a
semantic remainder, under coherence and universal totality. These results
are formalized in \texttt{Bridge/NonErasurePrinciple.lean} and
\texttt{Bridge/UniversalNonErasureLaw.lean}.

\paragraph{Infinity Compression (IC).}
The IC flagship~\cite{Spivack2026-IC-internal,Spivack2026-IC-summit,Spivack2026-IC-synthesis}
provides the internal origin of the program and the primary
discharge into the universal theory. Two distinct role-separated skeletons
with the same summit component extraction yield the witness-diversity data
for an adequate system. The IC adequate system forgets to the same RCS as
the IC universal instance ($\texttt{icAdequateReflectiveSystem\_toRCS} =
\texttt{rfl}$). The role permutation that swaps the two skeletons is a
non-identity $\compare$-preserving endomorphism, so IC also discharges
into the reflectively sufficient class.

\paragraph{NEMS / computability.}
The No External Model Selection (NEMS) framework~\cite{Spivack2026-NEMS01,Spivack2026-NEMS08,Spivack2026-NEMS81}
provides the outer (computability) layer. The diagonal barrier and semantic
remainder from NEMS instantiate the reflexive layer interface. NEMS
record-truth undecidability serves as the concrete semantic remainder.

\paragraph{APS / algebra.}
The abstract indexed programming systems (APS) framework~\cite{Spivack2026-APS}
provides the middle (realization) layer. The corrected exactness biconditional
(indexed composition $\Leftrightarrow$ finite tracking $+$ gluing) instantiates
the realization layer interface.

\paragraph{Cross-corpus synthesis.}
The three source programs combine into a concrete linked architecture
(\texttt{LinkedArchitecture}) with both coherence axioms proved definitionally.
The cross-corpus biconditional (IC enriched irreducibility
$\Leftrightarrow$ NEMS semantic remainder) and the universal non-erasure law
hold for all spine-eligible pairs. This synthesis is the concrete evidence that
the abstract residual geometry is not vacuous: it is instantiated by a real
cross-domain architecture with machine-checked proofs at every level.

\begin{remark}[Lean formalization]
\leavevmode\par\begin{itemize}[nosep,leftmargin=*,before=\raggedright]
\item \texttt{icAdequateReflectiveSystem} in \texttt{Summit/ICAdequateDischarge.lean}
\item \texttt{nemsReflexiveLayer} in \texttt{Instances/FromNEMS.lean}
\item \texttt{apsRealizationLayerFromIff} in \texttt{Instances/FromAPS.lean}
\item \texttt{crossCorpusLinkedArchitecture} in \texttt{Instances/CrossCorpusInstance.lean}
\item \texttt{universal\_nonerasure\_law} in \texttt{Bridge/UniversalNonErasureLaw.lean}
\end{itemize}
\end{remark}

%% ========================================================================
\section{Relationship to incompleteness and proof relevance}\label{sec:related}
%% ========================================================================

\paragraph{Relationship to Gödel's incompleteness.}
The fundamental equivalence has a structural parallel with Gödel's
incompleteness theorems~\cite{Godel1931}: in both cases, a system's own internal resources
produce the object that prevents its completeness. In Gödel's case, the
system's expressive power constructs an undecidable sentence; in ours, the
system's comparison geometry constructs a non-trivial fiber automorphism.
The all-or-nothing dichotomy (Theorem~\ref{thm:all-or-nothing}) mirrors
Gödel's result in that there is no partial escape: either the system is
completely exhaustive, or it has all five forms of non-exhaustion
simultaneously.

The differences are equally important. First, Gödel's theorem concerns
provability of sentences; ours concerns distinguishability of realizations.
The objects are different in kind. Second, Gödel's theorem is purely
negative; ours has a large positive component---the residual has a canonical
geometry with exact laws (the Fundamental Theorem of Residual Geometry,
Theorem~\ref{thm:fund-res}). Third, Gödel's theorem requires arithmetic
expressiveness; the fundamental theorem of residual geometry is
unconditional. The adequacy conditions that force nontriviality
(\S\ref{sec:adequacy}) are separate from the geometric theory itself.
The analogy is structural, not identificatory: the present theory is not a
new incompleteness theorem about provability, but a theorem about the
geometry of indistinguishability under maps.

\paragraph{The meta-level distinction.}
A natural question arises: if the residual can be measured exactly
(\S\ref{sec:optimal}), does this contradict Gödel's claim that a system
cannot know its own incompleteness?

No. The residual kernel is defined in terms of \emph{both} $\Bare$ and
$\Real$---it requires access to the realization level, which is exactly
what certification compresses away. Any predicate that factors through
$\compare$ (a ``bare-level'' predicate) is provably constant on every
kernel pair: it cannot detect whether the residual exists. Formally:
for any $Q : \Bare \to \mathrm{Prop}$ and any kernel pair $(x,y)$,
$Q(\compare(x)) = Q(\compare(y))$.

So the residual is \emph{invisible from the bare level} but
\emph{visible from the meta-level} (which has access to both $\Bare$ and
$\Real$). This is precisely the Gödelian structure: the veil exists and
has exact geometry, but the certification level cannot see it. An external
observer---or a stronger system with access to the realization level---can
measure the residual exactly. The certification itself cannot.

The theory does not let a system measure its own incompleteness. It lets
a \emph{meta-theory} measure the incompleteness of any system it can see
from outside. That is the correct reading of the resolution and measurement
results in \S\ref{sec:optimal}.

\paragraph{Related mathematical structures.}
The fiber decomposition of a map is classical; the residual kernel is the
complement of the diagonal in the fiber product $\Real \times_\Bare \Real$~\cite{MacLane1971}.
The observable algebra is a quotient Boolean algebra in the sense of
Stone duality~\cite{Stone1936}. The coloring interpretation connects to the classical theory
of chromatic numbers of complete graphs~\cite{BrooksBerge}. The functoriality results are
elementary category theory~\cite{MacLane1971}. What is new is not any single construction but
the \emph{synthesis}: the complete package of kernel, observables, refinement,
coloring, and complexity as a unified theory of information loss under arbitrary
maps, with the Classification Theorem and all-or-nothing dichotomy as the
governing structural results.

\paragraph{Why this is not dependent-type proof relevance.}
In dependent type theory~\cite{MartinLof1984,HoTT2013}, ``proof relevance'' means that different proofs of
the same proposition can be distinguished. Our non-exhaustion is superficially
similar: different realizations of the same bare certificate are distinguished.
But the resemblance is shallow. Proof relevance is a design choice in a type
system; our non-exhaustion is a \emph{theorem} about comparison maps, with
no dependence on any particular foundational framework. The residual kernel,
observable duality, and resolution theory have no analogues in the proof
relevance literature. The theory is about the geometry of what \emph{any}
certification process forgets, not about the internal structure of proofs.

%% ========================================================================
\section{The universal theory of forgetting}\label{sec:universal}
%% ========================================================================

The results of this paper are stated for an arbitrary comparison map
$\compare : \Real \to \Bare$ between arbitrary types. Nothing in the
definitions or theorems requires that $\compare$ arise from a ``certification
process.'' The word ``certification'' is a motivating interpretation, not a
formal hypothesis.

This means the entire theory applies to \emph{any function between any two
types}. Every such function determines a residual kernel, a fiber
decomposition, an observable algebra, a resolution theory, and exact
complexity bounds. The fundamental equivalence, the all-or-nothing dichotomy,
and the meta-level invisibility theorem hold universally.

\begin{theorem}[Universal forgetting]\label{thm:universal-forget}
For any function $\pi : E \to B$ between types in a common universe:
\begin{enumerate}[label=(\roman*),nosep]
\item The residual kernel
      $\{(x,y) \mid x \neq y \land \pi(x) = \pi(y)\}$ is the canonical
      object of what $\pi$ forgets.
\item Non-injectivity, witness diversity, fiber automorphism, and nontrivial
      kernel are equivalent (fundamental equivalence).
\item A predicate on $E$ factors through $\pi$ iff it does not separate any
      kernel pair (observable duality).
\item A refinement $r : E \to C$ makes $(\pi, r)$ injective iff $r$
      separates every kernel pair (resolution criterion).
\item Resolving a fiber of size $n$ requires $r$ to be injective on that
      fiber (resolution complexity).
\item Any predicate $Q : B \to \mathrm{Prop}$ satisfies
      $Q(\pi(x)) = Q(\pi(y))$ whenever $\pi(x) = \pi(y)$: the codomain
      level is blind to the residual (meta-level invisibility).
\item The full fundamental residual package and five-way equivalence exist
      unconditionally.
\end{enumerate}
\end{theorem}

This theorem applies simultaneously to:

\begin{center}
\begin{tabular}{lll}
\toprule
Domain & $\pi$ & Residual \\
\midrule
Group theory & homomorphism $\varphi : G \to H$ & classical kernel $\ker\varphi$ \\
Topology & quotient $X \to X/{\sim}$ & equivalence classes \\
Type theory & projection $\Sigma_{x:A} B(x) \to A$ & fiber family $B$ \\
Physics & measurement $\mathrm{state} \to \mathrm{observable}$ & unmeasured degrees of freedom \\
Compilation & compiler $\mathrm{source} \to \mathrm{bytecode}$ & erased source structure \\
Cryptography & hash $\mathrm{data} \to \mathrm{hash}$ & collision structure \\
\bottomrule
\end{tabular}
\end{center}

\noindent
In each case, the residual kernel is the set of distinct inputs that the
function identifies. The observable duality says which predicates can and
cannot see the lost information. The resolution complexity says how much
extra data is needed to recover it. And the meta-level theorem says the
lost information is invisible from the codomain.

The theory is therefore not merely a theory of certification failure. It is
a \emph{universal theory of information loss under any map}---a geometry of
forgetting that applies wherever a function compresses structure.

Section~\ref{sec:representation-as-structured-forgetting} sharpens this into a broader interpretive thesis---representation as selective forgetting and its reflexive consequences---explicitly layered on the formal spine (Remark~\ref{rem:representation-interpretive}).

\begin{remark}[Scope and non-claims]\label{rem:non-claims}
We do not claim that every useful map is best studied via this theory, nor
that all information-loss phenomena are exhausted by the kernel. We do not
claim automatic recoverability of the residual from the codomain: the
meta-level theorem (\S\ref{sec:related}) proves the opposite. We do not
claim that the theory replaces domain-specific tools (e.g., cohomological
obstructions in group extensions, or Rice-style barriers in computability).
The claim is that the residual kernel, observable duality, and resolution
theory provide a \emph{universal structural layer} that applies to any map,
complementing rather than replacing domain-specific analysis
(e.g.,~\cite{Brown1982} for group cohomology, \cite{Rice1953} for computability barriers).
\end{remark}

\begin{remark}[Lean formalization]
\leavevmode\par\begin{itemize}[nosep,leftmargin=*,before=\raggedright]
\item \texttt{rcsOfMap}, \texttt{forgettingKernel}, \texttt{universal\_fundamental\_\allowbreak{}equivalence}, \texttt{universal\_observable\_\allowbreak{}duality} in \texttt{Residual/\allowbreak{}UniversalForgetting.lean}
\item \texttt{universal\_resolution}, \texttt{universal\_meta\_blindness}, \texttt{every\_function\_has\_\allowbreak{}residual\_geometry}, \texttt{every\_function\_has\_five\_way\_equivalence} in the same file
\end{itemize}
\end{remark}

%% ========================================================================
\section{Representation as Structured Forgetting}
\label{sec:representation-as-structured-forgetting}
%% ========================================================================

\begin{remark}[Interpretive scope]\label{rem:representation-interpretive}
This section is a \emph{philosophical synthesis} layered on the machine-checked
geometry developed above. The elementary claims about fibers, collisions
$f(a_1)=f(a_2)$, and irrecoverability from the image alone restate standard
set-theoretic facts aligned with \S\S\ref{sec:kernel}--\ref{sec:refinement}.
In \emph{reflexive-closure-lean}, adjoining definitions package single-map
forgetting (\texttt{ForgetsDistinction}, non-injectivity), minimal
representation/remainder frames, class-level invariants
(\texttt{InvariantUnder}, \texttt{ClassInvariant}, \texttt{ConservationLaw}),
and the Paper~57 non-exhaustion bridge files named in Remark~\ref{rem:structured-forgetting-lean}.
Ascending claims about ``the function'' and ``ontology'' as selective regimes,
and the reflexive consequences drawn in \S\ref{sec:representation-as-structured-forgetting:reflexive},
remain interpretive where they go beyond those abstract interfaces: they
articulate a disciplined reading consistent with the formal package, not
additional theorems in \texttt{reflexive-architecture-lean}.
\end{remark}

The present analysis can now be sharpened into a more general claim. What maps forget is not an isolated pathology of a few badly behaved functions. It is the visible surface of a deeper structural fact: \emph{representation itself is selective}, and selectivity is already a form of forgetting. A representation is not a neutral duplication of what it represents. It is a constrained articulation of it. To articulate is to choose; to choose is to privilege some distinctions over others; and to privilege some distinctions is to leave others behind. Thus forgetting is not merely an occasional defect of representation. It is one of its constitutive conditions.

This point can be stated at several levels. At the most elementary level, a non-injective map forgets because distinct source points are identified in the image. At a higher level, the function that generates or governs such a map does not remove this loss, because the function is itself another structured semantic artifact, encoded under assumptions, constraints, and an interpretive stance. At a still higher level, ontology does not terminate the chain either, because an ontology is itself a selective articulation: a system of admissible categories, distinctions, and commitments. The result is a hierarchy in which one repeatedly ascends in search of closure, only to find that each new level is itself another regime of partial articulation. Representation therefore forms not only a chain of articulation, but a chain of forgetting.

\begin{keybox}{Representation is not possible without forgetting}
\textbf{Representation is not possible without forgetting.} To represent is not to copy a reality into a second place without loss. It is to impose a structured regime of admissible distinctions on it. Every such regime preserves some relations, suppresses others, and leaves a residue outside what it can fully articulate. Forgetting is therefore not merely something representation sometimes does. It is one of the structural conditions under which representation becomes possible at all.
\end{keybox}

\subsection{Maps Forget by Collapsing Distinctions}

Let \(f : A \to X\). Every point of \(A\) may indeed be sent to some point of \(X\). But this does not mean that nothing has been forgotten. What matters is not merely that each input receives an image, but whether distinct inputs remain distinguishable in the codomain. If there exist \(a_1 \neq a_2\) in \(A\) such that
\[
f(a_1)=f(a_2),
\]
then the map forgets the distinction between \(a_1\) and \(a_2\). The lost structure is not annihilated in some absolute sense; rather, it is rendered irrecoverable from the image alone. The codomain remembers only that the input lay somewhere in the fiber
\[
f^{-1}(x)=\{a\in A : f(a)=x\},
\]
not which member of that fiber it was.

This is the elementary geometry of forgetting. A map forgets not by failing to assign outputs, but by collapsing distinctions among inputs. Its fibers, kernel relations, and witness multiplicities are therefore not accidental byproducts. They are the formal anatomy of what the map cannot preserve. The map becomes a selective instrument of visibility: it lets some structure pass through while forcing other structure into equivalence classes.

\subsection{Conservation as class-level memory}
\label{sec:representation-conservation}

The dual question is not only ``what does this map forget?'' but ``what survives every transformation in an admissible family?''  At map level, non-injectivity forgets distinctions that are not invariant under passage through~$f$.  If one enlarges the viewpoint from a single map to a whole class~$\mathcal{F}$ of admissible maps $f \colon A \to A$ (time evolutions, symmetries, coarse-grainings allowed by a theory), a \emph{conservation law} can be read as a quantity or structure~$I$ such that for all $f \in \mathcal{F}$,
\[
  I(f(a)) = I(a),
\]
equivalently $I \circ f = I$ in the relevant functorial sense.  What is conserved is then exactly what no $f$ in the admissible class can erase from the perspective of~$I$.

\begin{keybox}{Conservation is class-level memory}
\textbf{A conservation law is a class-level memory.}  Instead of tracking the fiber-wise \emph{lost} remainder of one map, one tracks the \emph{unforgotten} remainder at family level: the structured core that persists under every admissible transformation.  Local forgetting asks what collapses under a single quotient; conservation asks what cannot be lost under the whole admissible regime.
\end{keybox}

This is not the same object as a single-map fiber remainder: class-level invariance picks out what is \emph{blind} to whichever distinctions each admissible map may merge---formally, any $I$ with $I \circ f = I$ is constant on fibers of~$f$.  The elegant duality is therefore
\begin{center}
\textit{Forgetting (one map): what is lost under $f$?}
\qquad
\textit{Conservation (a family): what can never be lost under every $f \in \mathcal{F}$?}
\end{center}
Physically, one may read the same structure twice: a dynamical law as a \emph{generator} of allowed maps, and simultaneously as a \emph{guarantee} that those maps preserve a common invariant core---a deeper notion of conservation than ``a number stays constant.''

The hierarchy suggested for the wider architecture is then
\[
  \text{single-map forgetting}
  \;\to\;
  \text{family-level invariance}
  \;\to\;
  \text{articulative non-exhaustion}
\]
local fibers and quotients on the left, conserved observables at the center, and (under the reflexive-closure stack) the persistence of semantic remainder under internal ascent on the right---so local and global pictures fold into one storyline rather than living in separate chapters.

Yet this first level already invites an objection. One may say: perhaps the image forgets, but the map itself does not. Perhaps everything apparently lost in the codomain is still preserved in the graph of the function. Perhaps the codomain forgets, but the total map remembers.

\subsection{The Function Does Not Close the Loss}

This reply fails because it treats the function as though it stood outside the economy of articulation in which the forgetting arose. But the function is not an external completion of the situation. It is itself another structured object. A function is not merely a bare assignment of outputs to inputs. It is an encoded law of transformation under a specific domain, a set of admissibility conditions, a constrained family of allowable operations, and an interpretive stance that determines what counts as an input, what counts as an output, and what structure is to matter in passing from one to the other.

In this sense a function is itself a bundle of meaning. It is not only a transformer of semantic structure; it is a semantic commitment concerning which distinctions are relevant, which relations are visible, which outputs count as legitimate, and which patterns are worth preserving. A function therefore carries a stance toward its domain. It says, in effect, both \emph{what to do} and \emph{what counts}. The very choice of codomain, the admissible input class, the invariants preserved, the quotienting tolerated, and the constraints imposed are all part of the semantic burden encoded by the function itself.

But then the appeal from the map to the function does not escape forgetting. It merely ascends to another structured articulation. If the codomain forgets by collapsing distinctions among inputs, the function forgets by encoding only one law of transformation under one regime of assumptions rather than another. It therefore cannot be treated as the final repository of everything the map left out. Its own articulation presupposes a background of unchosen possibilities, excluded distinctions, and unexpressed interpretive conditions. The function is thus not a neutral container of total meaning. It is another selective compression of it.

\begin{keybox}{The function does not recover what the map forgot}
\textbf{The function does not recover what the map forgot.} A function is itself a structured semantic artifact: a law of transformation encoded under assumptions, constraints, and an interpretive stance. To appeal from the image to the function is therefore not to escape remainder, but only to move one level upward within it. The higher-order articulation is itself another selective regime, and therefore another site of forgetting.
\end{keybox}

This point is decisive. If the image forgets and the function does not close that forgetting, then the chain does not end with the function. One may try to ascend again and say: perhaps the ontology within which the function is defined closes the gap. Perhaps the function is incomplete, but the ontology that determines its admissible objects and relations provides the final frame.

\subsection{Ontology Does Not Exhaust Meaning}

This move also fails. An ontology is not pure self-sufficiency. It is itself a structured selection: a choice of what kinds of entities, relations, categories, and admissible distinctions count within a given frame. An ontology says what there is, but in doing so it also says what sort of articulation will count as legitimate, what forms of difference are basic, what forms are derivative, and what kinds of explanation may begin. It is therefore not the elimination of choice, but its highest codified form within a given regime.

For this reason, ontology too belongs to the chain of selective articulation. It does not float above forgetting. It institutes forgetting at a deeper level. To choose an ontology is to decide which distinctions will be primary, which secondary, which invisible, and which unavailable. Even where the ontological commitments are taken as axiomatic, the problem is not removed. Axioms can terminate a derivation, but they do not terminate presupposition. An ontology stated as primitive still presupposes a basis on which its categories can be meaningful, its admissible distinctions can be drawn, and its commitments can be determinate as \emph{this} ontology rather than another.

Accordingly, ontology does not exhaust meaning. It is itself one more articulation within a wider semantic condition that it does not fully contain. The point here is not to specify that wider condition in positive metaphysical detail. The present argument requires only the negative conclusion: no ontology is self-grounding merely by being articulated as an ontology. If ontology is a regime of admissible choices, then the very fact that such choices can be meaningful already points beyond the ontology to what its own articulation presupposes.

\begin{keybox}{Ontology is selective articulation}
\textbf{Ontology is itself a selective articulation, and therefore a higher-order forgetting.} An ontology does not terminate remainder by naming the ultimate inventory of what there is. It is itself a structured choice of admissible categories, distinctions, and commitments. Even if such commitments are taken as axiomatic, their meaningfulness is not self-grounding. Ontology therefore does not exhaust the condition of its own articulation.
\end{keybox}

\subsection{The Chain of Representation Is a Chain of Forgetting}

The preceding steps may now be assembled into a single line of reasoning.

A map is a representation. It renders a source domain visible in a codomain by carrying certain distinctions forward and collapsing others. A function is a representation of a higher order. It encodes a law of transformation and a regime of admissibility, but in doing so it too selects and excludes. An ontology is a higher-order representation again. It codifies what sorts of objects, relations, and differences count at all. But it does so only by establishing one structured frame among others and therefore by instituting a deeper selectivity.

The chain thus has the following form:
\[
\text{map} \;\longrightarrow\; \text{function} \;\longrightarrow\; \text{ontology} \;\longrightarrow\; \cdots
\]
At each stage one may hope that the higher level recovers what the lower level forgot. But each higher level is itself another articulation, and articulation is possible only through selectivity. For that reason the chain is simultaneously
\[
\text{representation} \;\longrightarrow\; \text{higher-order representation} \;\longrightarrow\; \text{higher-order representation} \;\longrightarrow\; \cdots
\]
and
\[
\text{forgetting} \;\longrightarrow\; \text{higher-order forgetting} \;\longrightarrow\; \text{higher-order forgetting} \;\longrightarrow\; \cdots .
\]
One does not climb out of forgetting by ascending within the same economy of structured articulation. One only reiterates it in a more rarefied form.

This yields a deeper principle:

\begin{keybox}{Every internal articulation forgets}
\textbf{Every internal articulation is also an internal forgetting.} A map forgets by collapsing distinctions. A function forgets by encoding only one law of transformation under one regime of assumptions. An ontology forgets by selecting one admissible basis of categories and commitments rather than another. The ascent from map to function to ontology is therefore not a ladder to total recovery, but a hierarchy of increasingly general selective regimes.
\end{keybox}

This is why representation is impossible without forgetting. A representation that forgot nothing would preserve every distinction, every relation, every admissible interpretation, every contextual condition, and every possible articulation simultaneously. But such a structure would no longer function as a representation in any usable sense. It would cease to select, cease to compress, cease to interpret, and therefore cease to be representational at all. Representation requires visibility under constraint, and visibility under constraint is inseparable from exclusion.

\subsection{Where the Chain Leads}

The natural thought is that, after enough ascent, one should eventually arrive at a final level that contains everything the lower levels left out. The foregoing analysis shows why this hope must fail. Since each level is itself another selective articulation, the chain does not culminate in a final all-inclusive representation. It culminates instead in the recognition of a limit.

This limit should not be misunderstood as one last hidden item waiting to be added to the inventory. Nor is it merely one final forgotten content still lurking behind the current articulation. The point is deeper. The chain leads not to one more representable thing, but to what no such chain can fully contain. Every level of articulation can leave remainder because every level is itself constituted through selective admissibility. The final remainder is therefore not just another omitted object within the chain. It is what cannot be reduced to one more internally articulated item of the same kind.

For that reason, the deepest remainder is not properly described as forgotten content. Forgetting presupposes that something was, at least in principle, available to the relevant regime of knowledge or articulation and then lost, collapsed, omitted, or rendered irrecoverable within it. But the present chain points to something stronger: a limit of articulability itself. What lies at the end of the chain is not merely what a given map forgot, nor what a given function failed to encode, nor what a given ontology omitted from its inventory. It is what no internal articulation of this type can fully exhaust.

\begin{keybox}{Limit of articulation}
\textbf{The chain of forgetting terminates not in one final forgotten item, but at the limit of articulation itself.} What cannot be fully admitted into the regime of articulated knowing is not forgotten in the ordinary sense, because forgetting already presupposes prior admissibility to knowledge. The final remainder is deeper: it is irreducible non-exhaustion.
\end{keybox}

This gives the geometry of forgetting its widest significance. What maps forget is not merely a local fact about non-injective comparisons. It is the first visible manifestation of a general principle: every regime of internal articulation preserves some structure only by excluding other structure, and no ascent within such regimes closes that exclusion once and for all. The geometry of fibers, kernels, witness multiplicities, and exact resolution cost is therefore the opening layer of a more general architecture of non-exhaustion. The local remainder of a map, the higher-order remainder of a function, and the deepest remainder of ontology are not unrelated phenomena. They are successive expressions of one structural fact: closure is achievable only as partial articulation, and partial articulation never coincides with total semantic reality.

\subsection{Reflexive Consequence}
\label{sec:representation-as-structured-forgetting:reflexive}

The importance of this chain becomes greatest in the reflexive case. A reflexive system is one that does not merely contain maps, functions, and ontological commitments as inert contents, but is itself structured through transformations that return upon its own conditions of articulation. In such a setting the foregoing argument does not remain external to the system. It is inherited by the system from within. If representation is structurally inseparable from forgetting, and if reflexive articulation is still articulation, then a reflexive system cannot be expected to close itself into a final, exhaustive internal image. It may return to itself; it may partially articulate itself; it may even generate highly disciplined internal representations of its own structure. But each such articulation remains one more selective regime. Hence reflexive closure does not yield reflexive exhaustion.

The consequence is not merely that reflexive systems happen to be difficult to complete. It is that their incompletion is structurally principled. Their internal images, no matter how refined, remain subject to the same architecture that governs maps, functions, and ontologies generally. What appears in the local case as non-injective forgetting appears in the reflexive case as irreducible remainder. The geometry of what maps forget is thus not peripheral to reflexive limitation. It is one of its most elementary and intuitive expressions.

\begin{keybox}{Reflexive systems inherit the chain}
\textbf{Reflexive systems inherit the chain of forgetting from within.} Since a reflexive system articulates itself only through internal maps, functions, and ontological structures, it cannot coincide with a final exhaustive internal image of itself. Reflexive closure is therefore possible, but reflexive exhaustion is not. What maps forget in the local case becomes, at the reflexive level, a principled irreducible remainder.
\end{keybox}

The upshot is that the geometry of what maps forget should not be read merely as a theory of information loss under non-injective comparison. It is also a theory of why representation has a constitutive limit, why ascent within structured articulation does not abolish remainder, and why reflexive systems cannot eliminate the very distance through which they become representable to themselves in the first place. The local geometry of forgetting opens directly onto a general architecture of non-exhaustion.

\begin{remark}[Lean formalization: reflexive-closure bridge]\label{rem:structured-forgetting-lean}
Summit~II (\S\ref{sec:representation-as-structured-forgetting}) cites \emph{reflexive-closure-lean} for a thin, \texttt{sorry}-free layer that does not duplicate the universal residual formalization in this repository:
\leavevmode\par\begin{itemize}[nosep,leftmargin=*,before=\raggedright]
\item \texttt{ReflexiveClosure/\allowbreak{}Core/\allowbreak{}MapForgetting.lean}: \texttt{ForgetsDistinction}, \texttt{HasForgetting}, \texttt{hasForgetting\_iff\_not\_injective}, \texttt{representation\_forgetting\_theorem}.
\item \texttt{ReflexiveClosure/\allowbreak{}Core/\allowbreak{}FamilyConservation.lean}: \texttt{InvariantUnder} ($I \circ f = I$), \texttt{ClassInvariant}/\texttt{ConservationLaw}, \texttt{invariant\_agrees\_on\_fiber}, \texttt{invariant\_identifies\_forgotten\_distinctions}, \texttt{class\_level\_memory\_vs\_local\_forgetting}.
\item \texttt{ReflexiveClosure/\allowbreak{}Core/\allowbreak{}RepresentationFrame.lean}: \texttt{RepSemanticRemainder}, \texttt{RepresentationExhaustive}, \texttt{RepresentationSelective}.
\item \texttt{ReflexiveUnfolding/\allowbreak{}Bridge/\allowbreak{}StructuredForgetting.lean}: alignment with Paper~57 \texttt{SemanticRemainder}.
\item \texttt{ReflexiveUnfolding/\allowbreak{}Theorems/\allowbreak{}FunctionalNonExhaustion.lean}, \texttt{ArticulativeAscent.lean}: functional and ascent non-exhaustion downstream of Papers~56--57.
\end{itemize}
\textbf{Trust boundary:} these are abstract type-theoretic interfaces; conservation is formalized for endomorphisms $\alpha \to \alpha$. Interpreting $\mathcal{F}$ as ``physical evolution'' or ``fundamental symmetry'' is not itself machine-checked.
\end{remark}

%% ========================================================================
\section{Frontier}\label{sec:frontier}
%% ========================================================================

The theory as presented is complete at the qualitative and structural level.
We distinguish immediate mathematical next steps from deeper programmatic
conjectures.

\subsubsection*{Immediate next steps}

\paragraph{Residual spectra and moduli.}
The fiber-size spectrum (the multiset of fiber sizes) is a complete invariant
of resolution complexity. Defining entropy-like measures of residual
spread, moduli of kernel structure, and spectral invariants that survive
bounded refinement would turn the theory from structural to quantitative.

\paragraph{Residual observables under refinement.}
Which non-bare-determined predicates survive a given refinement? Can one
classify the ``persistent residual observables'' that remain non-bare-determined
after any bounded refinement? This is the formal version of asking what
structure is \emph{robustly} invisible to certification.

\paragraph{Iterated certification.}
What happens when a refined comparison is itself treated as a new comparison
map and the theory is applied again? Does the residual stabilize, shrink
monotonically, or exhibit more complex dynamics?

\paragraph{Category-theoretic packaging.}
Compatible maps compose associatively with identity units
(Theorem~\ref{thm:functorial}). The next step is to package RCS and
compatible maps into a category and ask whether the kernel construction
defines a functor, and whether the fundamental residual package is natural.

\subsubsection*{Deeper conjectures}

\paragraph{Domain-specific kernel structure.}
In each discharge domain (group extensions, embedding problems, Quillen
fibers, computability barriers), the kernel has domain-specific structure
beyond the abstract theory. Rewriting these domain results in kernel
language would test the theory's power as a unifying instrument.

\paragraph{Broader inevitability.}
The current summit classes require explicit structural data. Whether there
exist broader ``richness'' conditions that force witness diversity or fiber
automorphisms from more abstract principles---without supplying witnesses
or twists as data---remains the deepest open question.

%% ========================================================================
\section{Conclusion}\label{sec:conclusion}
%% ========================================================================

This paper has established that information loss under any map is not a
deficiency to be lamented but a geometric phenomenon to be studied. The
residual kernel is the canonical object of that geometry; the fundamental
equivalence theorem shows that non-injectivity, witness diversity, fiber
automorphism, and nontrivial kernel are the same phenomenon viewed from
five angles; and the fundamental theorem of residual geometry gives the
complete structure of the residual unconditionally for every function
between any two types.

The theory opens a new kind of mathematical question. Classical algebra
asks: given a group, what are its subgroups? Classical topology asks: given
a space, what are its open sets? The residual theory asks: given \emph{any
function}, what is the exact geometry of what it forgets? The answer---a
disjoint union of complete indistinguishability cliques, with observables
dual to kernel separation and resolution governed by graph coloring---is
both precise and universal.

The formalization in Lean~4 with zero \texttt{sorry} ensures that every
claim in this paper is machine-checked. The discharge against the Infinity
Compression flagship, the NEMS computability anchor, and the APS algebraic
layer ensures that the abstract theory is not vacuous.

Section~\ref{sec:representation-as-structured-forgetting} connects the fiber
geometry to a wider interpretive thesis: representation requires selective
articulation; appealing from image to function to ontology reiterates rather
than closes remainder; dual to single-map forgetting, class-level invariants
(\S\ref{sec:representation-conservation}) express what an admissible family
of transformations can never erase; and reflexive systems inherit that
architecture internally (Remark~\ref{rem:representation-interpretive} for
scope; Remark~\ref{rem:structured-forgetting-lean} for the small formal
bridge). The mathematical core remains Theorem~\ref{thm:fund-res} and the
universal package---closure of articulation without pretense of total internal
exhaustion is philosophy consistent with, not identical to, those lemmas.

What a non-injective map forgets is not an accident, not an informal remainder,
and not merely a by-product of examples. It is a canonically organized object
with exact laws. This paper identifies that object, proves its classification,
determines its observables, and computes the exact cost of resolving it. The
problem is no longer whether information is lost under a map---it always is,
when the map is non-injective. The problem is to understand the geometry of
that loss. This paper gives that geometry.

%% ========================================================================
\appendix
\section{Theorem index and Lean correspondence}\label{app:index}
%% ========================================================================

{\sloppy The table below lists every named paper result, its Lean theorem name, and
its module. All names are within the \texttt{ReflexiveArchitecture\allowbreak.Universal}
namespace; \texttt{Summit/*} live under
\texttt{ReflexiveArchitecture\allowbreak.Universal\allowbreak.Summit} and
\texttt{Residual/*} under
\texttt{ReflexiveArchitecture\allowbreak.Universal\allowbreak.Residual}.}

\begingroup\sloppy
\begin{longtable}{>{\footnotesize\raggedright\arraybackslash}p{4.4cm} >{\footnotesize\raggedright\arraybackslash}p{6.4cm} >{\footnotesize\raggedright\arraybackslash}p{3.8cm}}
\toprule
\footnotesize\textbf{Paper result} & \footnotesize\textbf{Lean name} & \footnotesize\textbf{Module} \\
\midrule
\endhead
\midrule
\multicolumn{3}{r}{\small(continued)}\\
\endfoot
\bottomrule
\endlastfoot
%
Classification Theorem (Thm.~\ref{thm:fund-equiv}) &
  \texttt{equiv\_1\_iff\_2} \ldots \texttt{equiv\_1\_iff\_5} &
  \texttt{Summit/\allowbreak{}FundamentalEquivalence} \\[4pt]
%
All-or-nothing dichotomy (Thm.~\ref{thm:all-or-nothing}) &
  \texttt{all\_or\_nothing} &
  \texttt{Summit/\allowbreak{}FundamentalEquivalence} \\[4pt]
%
Five-way equivalence package &
  \texttt{FiveWayEquivalence}, \texttt{fiveWayEquivalence} &
  \texttt{Summit/\allowbreak{}FundamentalEquivalence} \\[4pt]
%
Fiber automorphism diagonal &
  \texttt{nonExhaustive\_implies\_\allowbreak{}fiberAutomorphism} &
  \texttt{Summit/\allowbreak{}FundamentalEquivalence} \\[4pt]
%
Vanishing Criterion (Thm.~\ref{thm:exhaust}) &
  \texttt{exhaustive\_iff\_\allowbreak{}kernel\_empty},
  \texttt{kernel\_empty\_iff\_\allowbreak{}all\_fibers\_\allowbreak{}subsingleton} &
  \texttt{Residual/\allowbreak{}FundamentalTheorem} \\[4pt]
%
Adequacy Theorem (Thm.~\ref{thm:adequate-ne}) &
  \texttt{adequate\_nonExhaustive} &
  \texttt{Summit/\allowbreak{}Adequacy} \\[4pt]
%
Self-generation (Thm.~\ref{thm:self-gen}) &
  \texttt{reflectivelySufficient\_\allowbreak{}nonExhaustive},
  \texttt{toAdequate} &
  \texttt{Summit/\allowbreak{}SelfGeneration} \\[4pt]
%
Kernel characterization (Thm.~\ref{thm:kernel-char}) &
  \texttt{residualKernel\_\allowbreak{}nonempty\_iff},
  \texttt{residualKernel\_eq\_\allowbreak{}of\_compare\_eq} &
  \texttt{Residual/\allowbreak{}ResidualKernel} \\[4pt]
%
Fiberwise decomposition (Thm.~\ref{thm:fiber-decomp}) &
  \texttt{residualKernel\_eq\_\allowbreak{}iUnion\_kernelAt} &
  \texttt{Residual/\allowbreak{}KernelStratification} \\[4pt]
%
Complete-fiber theorem (Thm.~\ref{thm:complete-fiber}) &
  \texttt{kernelAt\_complete} &
  \texttt{Residual/\allowbreak{}MinimalResolution} \\[4pt]
%
Connectivity theorem (Thm.~\ref{thm:connectivity}) &
  \texttt{kernelRelRefl\_equivalence},
  \texttt{fiber\_eq\_\allowbreak{}kernelRelRefl\_class} &
  \texttt{Residual/\allowbreak{}KernelGraph} \\[4pt]
%
Observable algebra (Thm.~\ref{thm:obs-algebra}) &
  \texttt{bareDetermined\_and}, \texttt{bareDetermined\_or},
  \texttt{bareDetermined\_not}, \ldots &
  \texttt{Residual/\allowbreak{}ObservableAlgebra} \\[4pt]
%
Universal factorization (Thm.~\ref{thm:univ-factor}) &
  \texttt{constant\_on\_fibers\_\allowbreak{}iff\_factors} &
  \texttt{Residual/\allowbreak{}ObservableAlgebra} \\[4pt]
%
Kernel-observable duality (Thm.~\ref{thm:duality}) &
  \texttt{not\_bareDetermined\_\allowbreak{}iff\_separates\_kernel} &
  \texttt{Residual/\allowbreak{}ResidualKernel} \\[4pt]
%
Resolution criterion (Thm.~\ref{thm:resolution}) &
  \texttt{resolves\_iff\_separates\_\allowbreak{}all\_kernel\_pairs} &
  \texttt{Residual/\allowbreak{}ResidualKernel} \\[4pt]
%
Resolution complexity (Thm.~\ref{thm:res-complex}) &
  \texttt{resolution\_complexity\_\allowbreak{}theorem} &
  \texttt{Residual/\allowbreak{}FundamentalTheorem} \\[4pt]
%
Certification--Refinement Gap (Thm.~\ref{thm:cert-gap}) &
  \texttt{certification\_refinement\_gap} &
  \texttt{Residual/\allowbreak{}OptimalCertification} \\[4pt]
%
Residual Geometry Theorem (Thm.~\ref{thm:fund-res}) &
  \texttt{FundamentalResidualPackage},
  \texttt{fundamentalResidualPackage} &
  \texttt{Residual/\allowbreak{}FundamentalTheorem} \\[4pt]
%
Persistence dichotomy (Thm.~\ref{thm:persist}) &
  \texttt{persistence\_dichotomy},
  \texttt{exists\_persistent\_\allowbreak{}of\_unresolved} &
  \texttt{Residual/\allowbreak{}MetaLevel} \\[4pt]
%
Universal forgetting (Thm.~\ref{thm:universal-forget}) &
  \texttt{every\_function\_has\_\allowbreak{}residual\_geometry},
  \texttt{every\_function\_has\_\allowbreak{}five\_way\_equivalence} &
  \texttt{Residual/\allowbreak{}UniversalForgetting} \\[4pt]
%
Nothing remains at optimal (Thm.~\ref{thm:nothing-remains}) &
  \texttt{nothing\_remains\_at\_optimal} &
  \texttt{Residual/\allowbreak{}OptimalCertification} \\[4pt]
%
Category laws (Thm.~\ref{thm:functorial}) &
  \texttt{RCSCategoryLaws}, \texttt{rcsCategoryLaws} &
  \texttt{Residual/\allowbreak{}RCSCategory} \\
\end{longtable}
\endgroup

\noindent
Build: \texttt{lake build} from \texttt{reflexive-architecture-lean/}. The table above is the
paper-facing index; as of this revision the summit route comprises on the order of $235$ theorems
across $26$ modules under \texttt{ReflexiveArchitecture.Universal} (residual and summit classes).
Supplementary bridge files in \texttt{reflexive-closure-lean} appear in Remark~\ref{rem:structured-forgetting-lean}.

%% ========================================================================
\begin{thebibliography}{99}
%% ========================================================================

%%% Standard canonical references

\bibitem{Godel1931}
K.~Gödel,
\emph{Über formal unentscheidbare Sätze der Principia Mathematica und
verwandter Systeme~I},
Monatshefte für Mathematik und Physik, 38:173--198, 1931.
English translation: \emph{On Formally Undecidable Propositions of
Principia Mathematica and Related Systems}, Basic Books, 1962.

\bibitem{Stone1936}
M.~H.~Stone,
\emph{The theory of representations for Boolean algebras},
Transactions of the American Mathematical Society, 40(1):37--111, 1936.

\bibitem{BrooksBerge}
R.~L.~Brooks,
\emph{On colouring the nodes of a network},
Mathematical Proceedings of the Cambridge Philosophical Society,
37(2):194--197, 1941;
C.~Berge,
\emph{Graphs and Hypergraphs},
North-Holland, Amsterdam, 1973.

\bibitem{MacLane1971}
S.~Mac~Lane,
\emph{Categories for the Working Mathematician},
Springer-Verlag, New York, 1971.

\bibitem{Rice1953}
H.~G.~Rice,
\emph{Classes of recursively enumerable sets and their decision problems},
Transactions of the American Mathematical Society, 74(2):358--366, 1953.

\bibitem{MartinLof1984}
P.~Martin-Löf,
\emph{Intuitionistic Type Theory},
Bibliopolis, Naples, 1984.

\bibitem{HoTT2013}
The Univalent Foundations Program,
\emph{Homotopy Type Theory: Univalent Foundations of Mathematics},
Institute for Advanced Study, 2013.
Available at \url{https://homotopytypetheory.org/book/}.

\bibitem{Lean4}
L.~de~Moura and S.~Ullrich,
\emph{The Lean~4 Theorem Prover and Programming Language},
in \emph{Automated Deduction --- CADE~28}, Lecture Notes in Computer Science
vol.~12699, Springer, 2021, pp.~625--635.

\bibitem{Mathlib}
The Mathlib Community,
\emph{The Lean Mathematical Library},
in \emph{Proceedings of CPP 2020}, ACM, 2020, pp.~367--381.

\bibitem{Brown1982}
K.~S.~Brown,
\emph{Cohomology of Groups},
Graduate Texts in Mathematics vol.~87, Springer-Verlag, New York, 1982.

%%% Source program references

\bibitem{Spivack2026-NEMS01}
N.~Spivack,
\emph{The NEMS Framework: No External Model Selection},
Paper~1, NEMS Suite, 2026.
Lean: \texttt{nems-lean}; zero \texttt{sorry}, zero custom axioms.

\bibitem{Spivack2026-NEMS08}
N.~Spivack,
\emph{From NEMS to MFRR: A Machine-Checked Bridge Between Semantic Closure
and Reflexive Reality},
Paper~8, NEMS Suite, 2026.
Lean: \texttt{nems-lean}; zero \texttt{sorry}, zero custom axioms (8051 jobs).

\bibitem{Spivack2026-NEMS15}
N.~Spivack,
\emph{No Emulation, Self-Necessitating Adjudication},
Paper~15, NEMS Suite, 2026.
Lean: \texttt{nems-lean}; zero \texttt{sorry}, zero custom axioms (8069 jobs).

\bibitem{Spivack2026-NEMS81}
N.~Spivack,
\emph{Big Results from the NEMS Program},
Paper~81, NEMS Suite (Synthesis), 2026.
Lean: \texttt{nems-lean}.

\bibitem{Spivack2026-APS}
N.~Spivack,
\emph{Composition as Finite Tracking plus Gluing, and its Separation from
Recursion in Abstract Indexed Programming Systems},
2026.
Lean: \texttt{aps-recursion-uniformization-lean}; zero \texttt{sorry},
zero additional axioms.

\bibitem{Spivack2026-IC-internal}
N.~Spivack,
\emph{Canonical Certification Does Not Exhaust Reflective Structure},
Infinity Compression Suite, 2026.
Lean: \texttt{infinity-compression-lean}; zero \texttt{sorry}.

\bibitem{Spivack2026-IC-synthesis}
N.~Spivack,
\emph{Certification, Realization, and Obstruction---A Universal Fiber Architecture},
Infinity Compression Suite (Synthesis Map), 2026.
Lean: \texttt{infinity-compression-lean}.

\bibitem{Spivack2026-IC-summit}
N.~Spivack,
\emph{Reflective Non-Exhaustion Summit},
Infinity Compression Suite, 2026.
Lean: \texttt{infinity-compression-lean}.

\bibitem{Spivack2026-Strata}
N.~Spivack,
\emph{Closure, Realization, and Reflective Residue},
Reflexive Architecture / Strata, 2026.
Lean: \texttt{reflexive-architecture-lean}; zero \texttt{sorry},
zero custom axioms ($\sim$8270 jobs).

\end{thebibliography}

\end{document}